% Môn: Vật lý
% Lớp: 10
% Chương: Chương V. Động lượng
% Bài: L10C5 Bài 30. Thực hành: Xác định động lượng của vật trước và sau va chạm
% Số câu: 162
% App đọc file này trực tiếp — sửa rồi Commit trên GitHub, bấm Đồng bộ GitHub .tex

% L10C5 Bài 30. Thực hành: Xác định động lượng của vật trước và sau va chạm

% ===== Câu 1 =====
% ID: AIQ_L10C5_FULL_0325
% Mức: NB
\dangbt{Nhận biết dụng cụ và bố trí thí nghiệm va chạm}
\begin{ex}
% ID: AIQ_L10C5_FULL_0325
% Mức: NB
Trong bố trí thí nghiệm va chạm có 2 cổng quang điện độc lập. Có bao nhiêu thiết bị trực tiếp ghi thời điểm xe đi qua?
\choice
{\True $2\ \mathrm{thiết\,bị}$}
{$3\ \mathrm{thiết\,bị}$}
{$4\ \mathrm{thiết\,bị}$}
{$1\ \mathrm{thiết\,bị}$}
\loigiai{
Mỗi cổng quang điện là một thiết bị ghi thời điểm, nên có 2 thiết bị. Suy ra kết quả $2\ \mathrm{thiết\,bị}$.
}
\end{ex}

% ===== Câu 2 =====
% ID: AIQ_L10C5_FULL_0326
% Mức: NB
\begin{ex}
% ID: AIQ_L10C5_FULL_0326
% Mức: NB
Trong bố trí thí nghiệm va chạm có 3 cổng quang điện độc lập. Có bao nhiêu thiết bị trực tiếp ghi thời điểm xe đi qua?
\choice
{$4\ \mathrm{thiết\,bị}$}
{\True $3\ \mathrm{thiết\,bị}$}
{$6\ \mathrm{thiết\,bị}$}
{$2\ \mathrm{thiết\,bị}$}
\loigiai{
Mỗi cổng quang điện là một thiết bị ghi thời điểm, nên có 3 thiết bị. Suy ra kết quả $3\ \mathrm{thiết\,bị}$.
}
\end{ex}

% ===== Câu 3 =====
% ID: AIQ_L10C5_FULL_0327
% Mức: TH
\begin{ex}
% ID: AIQ_L10C5_FULL_0327
% Mức: TH
Trong bố trí thí nghiệm va chạm có 4 cổng quang điện độc lập. Có bao nhiêu thiết bị trực tiếp ghi thời điểm xe đi qua?
\choice
{$5\ \mathrm{thiết\,bị}$}
{$8\ \mathrm{thiết\,bị}$}
{\True $4\ \mathrm{thiết\,bị}$}
{$3\ \mathrm{thiết\,bị}$}
\loigiai{
Mỗi cổng quang điện là một thiết bị ghi thời điểm, nên có 4 thiết bị. Suy ra kết quả $4\ \mathrm{thiết\,bị}$.
}
\end{ex}

% ===== Câu 4 =====
% ID: AIQ_L10C5_FULL_0328
% Mức: TH
\begin{ex}
% ID: AIQ_L10C5_FULL_0328
% Mức: TH
Trong bố trí thí nghiệm va chạm có 5 cổng quang điện độc lập. Có bao nhiêu thiết bị trực tiếp ghi thời điểm xe đi qua?
\choice
{$6\ \mathrm{thiết\,bị}$}
{$10\ \mathrm{thiết\,bị}$}
{$4\ \mathrm{thiết\,bị}$}
{\True $5\ \mathrm{thiết\,bị}$}
\loigiai{
Mỗi cổng quang điện là một thiết bị ghi thời điểm, nên có 5 thiết bị. Suy ra kết quả $5\ \mathrm{thiết\,bị}$.
}
\end{ex}

% ===== Câu 5 =====
% ID: AIQ_L10C5_FULL_0329
% Mức: VD
\begin{ex}
% ID: AIQ_L10C5_FULL_0329
% Mức: VD
Trong bố trí thí nghiệm va chạm có 2 cổng quang điện độc lập. Có bao nhiêu thiết bị trực tiếp ghi thời điểm xe đi qua?
\choice
{\True $2\ \mathrm{thiết\,bị}$}
{$3\ \mathrm{thiết\,bị}$}
{$4\ \mathrm{thiết\,bị}$}
{$1\ \mathrm{thiết\,bị}$}
\loigiai{
Mỗi cổng quang điện là một thiết bị ghi thời điểm, nên có 2 thiết bị. Suy ra kết quả $2\ \mathrm{thiết\,bị}$.
}
\end{ex}

% ===== Câu 6 =====
% ID: AIQ_L10C5_FULL_0330
% Mức: VD
\begin{ex}
% ID: AIQ_L10C5_FULL_0330
% Mức: VD
Trong bố trí thí nghiệm va chạm có 3 cổng quang điện độc lập. Có bao nhiêu thiết bị trực tiếp ghi thời điểm xe đi qua?
\choice
{$4\ \mathrm{thiết\,bị}$}
{\True $3\ \mathrm{thiết\,bị}$}
{$6\ \mathrm{thiết\,bị}$}
{$2\ \mathrm{thiết\,bị}$}
\loigiai{
Mỗi cổng quang điện là một thiết bị ghi thời điểm, nên có 3 thiết bị. Suy ra kết quả $3\ \mathrm{thiết\,bị}$.
}
\end{ex}

% ===== Câu 7 =====
% ID: AIQ_L10C5_FULL_0331
% Mức: VD
\begin{ex}
% ID: AIQ_L10C5_FULL_0331
% Mức: VD
Trong bố trí thí nghiệm va chạm có 4 cổng quang điện độc lập. Có bao nhiêu thiết bị trực tiếp ghi thời điểm xe đi qua?
\choice
{$5\ \mathrm{thiết\,bị}$}
{$8\ \mathrm{thiết\,bị}$}
{\True $4\ \mathrm{thiết\,bị}$}
{$3\ \mathrm{thiết\,bị}$}
\loigiai{
Mỗi cổng quang điện là một thiết bị ghi thời điểm, nên có 4 thiết bị. Suy ra kết quả $4\ \mathrm{thiết\,bị}$.
}
\end{ex}

% ===== Câu 8 =====
% ID: AIQ_L10C5_FULL_0332
% Mức: VDC
\begin{ex}
% ID: AIQ_L10C5_FULL_0332
% Mức: VDC
Trong bố trí thí nghiệm va chạm có 5 cổng quang điện độc lập. Có bao nhiêu thiết bị trực tiếp ghi thời điểm xe đi qua?
\choice
{$6\ \mathrm{thiết\,bị}$}
{$10\ \mathrm{thiết\,bị}$}
{$4\ \mathrm{thiết\,bị}$}
{\True $5\ \mathrm{thiết\,bị}$}
\loigiai{
Mỗi cổng quang điện là một thiết bị ghi thời điểm, nên có 5 thiết bị. Suy ra kết quả $5\ \mathrm{thiết\,bị}$.
}
\end{ex}

% ===== Câu 9 =====
% ID: AIQ_L10C5_FULL_0333
% Mức: VDC
\begin{ex}
% ID: AIQ_L10C5_FULL_0333
% Mức: VDC
Trong bố trí thí nghiệm va chạm có 2 cổng quang điện độc lập. Có bao nhiêu thiết bị trực tiếp ghi thời điểm xe đi qua?
\choice
{\True $2\ \mathrm{thiết\,bị}$}
{$3\ \mathrm{thiết\,bị}$}
{$4\ \mathrm{thiết\,bị}$}
{$1\ \mathrm{thiết\,bị}$}
\loigiai{
Mỗi cổng quang điện là một thiết bị ghi thời điểm, nên có 2 thiết bị. Suy ra kết quả $2\ \mathrm{thiết\,bị}$.
}
\end{ex}

% ===== Câu 10 =====
% ID: AIQ_L10C5_FULL_0334
% Mức: VDC
\begin{ex}
% ID: AIQ_L10C5_FULL_0334
% Mức: VDC
Trong bố trí thí nghiệm va chạm có 3 cổng quang điện độc lập. Có bao nhiêu thiết bị trực tiếp ghi thời điểm xe đi qua?
\choice
{$4\ \mathrm{thiết\,bị}$}
{\True $3\ \mathrm{thiết\,bị}$}
{$6\ \mathrm{thiết\,bị}$}
{$2\ \mathrm{thiết\,bị}$}
\loigiai{
Mỗi cổng quang điện là một thiết bị ghi thời điểm, nên có 3 thiết bị. Suy ra kết quả $3\ \mathrm{thiết\,bị}$.
}
\end{ex}

% ===== Câu 11 =====
% ID: AIQ_L10C5_FULL_0335
% Mức: NB
\begin{ex}
% ID: AIQ_L10C5_FULL_0335
% Mức: NB
Xét các phát biểu thuộc dạng “Nhận biết dụng cụ và bố trí thí nghiệm va chạm” ở tình huống số 1.
\choiceTF
{\True Cổng quang điện dùng đo thời gian.}
{\True Cân dùng đo khối lượng.}
{Máng nghiêng bắt buộc phải rung.}
{\True Xe phải chuyển động gần như thẳng.}
\loigiai{
\begin{itemchoice}
\item (Đúng) Cổng quang điện dùng đo thời gian. (Vì): Phù hợp với định nghĩa hoặc hệ thức vật lí. b) Đúng: Phù hợp với định nghĩa hoặc hệ thức vật lí. c) Sai: Không phù hợp với định nghĩa hoặc hệ thức vật lí. d) Đúng: Phù hợp với định nghĩa hoặc hệ thức vật lí
\item (Đúng) Cân dùng đo khối lượng.
\item (Sai) Máng nghiêng bắt buộc phải rung.
\item (Đúng) Xe phải chuyển động gần như thẳng.
\end{itemchoice}
}
\end{ex}

% ===== Câu 12 =====
% ID: AIQ_L10C5_FULL_0336
% Mức: TH
\begin{ex}
% ID: AIQ_L10C5_FULL_0336
% Mức: TH
Xét các phát biểu thuộc dạng “Nhận biết dụng cụ và bố trí thí nghiệm va chạm” ở tình huống số 2.
\choiceTF
{\True Cổng quang điện dùng đo thời gian.}
{\True Cân dùng đo khối lượng.}
{Máng nghiêng bắt buộc phải rung.}
{\True Xe phải chuyển động gần như thẳng.}
\loigiai{
\begin{itemchoice}
\item (Đúng) Cổng quang điện dùng đo thời gian. (Vì): Phù hợp với định nghĩa hoặc hệ thức vật lí. b) Đúng: Phù hợp với định nghĩa hoặc hệ thức vật lí. c) Sai: Không phù hợp với định nghĩa hoặc hệ thức vật lí. d) Đúng: Phù hợp với định nghĩa hoặc hệ thức vật lí
\item (Đúng) Cân dùng đo khối lượng.
\item (Sai) Máng nghiêng bắt buộc phải rung.
\item (Đúng) Xe phải chuyển động gần như thẳng.
\end{itemchoice}
}
\end{ex}

% ===== Câu 13 =====
% ID: AIQ_L10C5_FULL_0337
% Mức: TH
\begin{ex}
% ID: AIQ_L10C5_FULL_0337
% Mức: TH
Xét các phát biểu thuộc dạng “Nhận biết dụng cụ và bố trí thí nghiệm va chạm” ở tình huống số 3.
\choiceTF
{\True Cổng quang điện dùng đo thời gian.}
{\True Cân dùng đo khối lượng.}
{Máng nghiêng bắt buộc phải rung.}
{\True Xe phải chuyển động gần như thẳng.}
\loigiai{
\begin{itemchoice}
\item (Đúng) Cổng quang điện dùng đo thời gian. (Vì): Phù hợp với định nghĩa hoặc hệ thức vật lí. b) Đúng: Phù hợp với định nghĩa hoặc hệ thức vật lí. c) Sai: Không phù hợp với định nghĩa hoặc hệ thức vật lí. d) Đúng: Phù hợp với định nghĩa hoặc hệ thức vật lí
\item (Đúng) Cân dùng đo khối lượng.
\item (Sai) Máng nghiêng bắt buộc phải rung.
\item (Đúng) Xe phải chuyển động gần như thẳng.
\end{itemchoice}
}
\end{ex}

% ===== Câu 14 =====
% ID: AIQ_L10C5_FULL_0338
% Mức: VD
\begin{ex}
% ID: AIQ_L10C5_FULL_0338
% Mức: VD
Xét các phát biểu thuộc dạng “Nhận biết dụng cụ và bố trí thí nghiệm va chạm” ở tình huống số 4.
\choiceTF
{\True Cổng quang điện dùng đo thời gian.}
{\True Cân dùng đo khối lượng.}
{Máng nghiêng bắt buộc phải rung.}
{\True Xe phải chuyển động gần như thẳng.}
\loigiai{
\begin{itemchoice}
\item (Đúng) Cổng quang điện dùng đo thời gian. (Vì): Phù hợp với định nghĩa hoặc hệ thức vật lí. b) Đúng: Phù hợp với định nghĩa hoặc hệ thức vật lí. c) Sai: Không phù hợp với định nghĩa hoặc hệ thức vật lí. d) Đúng: Phù hợp với định nghĩa hoặc hệ thức vật lí
\item (Đúng) Cân dùng đo khối lượng.
\item (Sai) Máng nghiêng bắt buộc phải rung.
\item (Đúng) Xe phải chuyển động gần như thẳng.
\end{itemchoice}
}
\end{ex}

% ===== Câu 15 =====
% ID: AIQ_L10C5_FULL_0339
% Mức: VDC
\begin{ex}
% ID: AIQ_L10C5_FULL_0339
% Mức: VDC
Xét các phát biểu thuộc dạng “Nhận biết dụng cụ và bố trí thí nghiệm va chạm” ở tình huống số 5.
\choiceTF
{\True Cổng quang điện dùng đo thời gian.}
{\True Cân dùng đo khối lượng.}
{Máng nghiêng bắt buộc phải rung.}
{\True Xe phải chuyển động gần như thẳng.}
\loigiai{
\begin{itemchoice}
\item (Đúng) Cổng quang điện dùng đo thời gian. (Vì): Phù hợp với định nghĩa hoặc hệ thức vật lí. b) Đúng: Phù hợp với định nghĩa hoặc hệ thức vật lí. c) Sai: Không phù hợp với định nghĩa hoặc hệ thức vật lí. d) Đúng: Phù hợp với định nghĩa hoặc hệ thức vật lí
\item (Đúng) Cân dùng đo khối lượng.
\item (Sai) Máng nghiêng bắt buộc phải rung.
\item (Đúng) Xe phải chuyển động gần như thẳng.
\end{itemchoice}
}
\end{ex}

% ===== Câu 16 =====
% ID: AIQ_L10C5_FULL_0340
% Mức: TH
\begin{ex}
% ID: AIQ_L10C5_FULL_0340
% Mức: TH
Trong bố trí thí nghiệm va chạm có 4 cổng quang điện độc lập. Có bao nhiêu thiết bị trực tiếp ghi thời điểm xe đi qua?
\shortans{4}
\loigiai{
Mỗi cổng quang điện là một thiết bị ghi thời điểm, nên có 4 thiết bị. Suy ra kết quả $4\ \mathrm{thiết\,bị}$. Đáp án trả lời ngắn không ghi đơn vị.
}
\end{ex}

% ===== Câu 17 =====
% ID: AIQ_L10C5_FULL_0341
% Mức: TH
\begin{ex}
% ID: AIQ_L10C5_FULL_0341
% Mức: TH
Trong bố trí thí nghiệm va chạm có 5 cổng quang điện độc lập. Có bao nhiêu thiết bị trực tiếp ghi thời điểm xe đi qua?
\shortans{5}
\loigiai{
Mỗi cổng quang điện là một thiết bị ghi thời điểm, nên có 5 thiết bị. Suy ra kết quả $5\ \mathrm{thiết\,bị}$. Đáp án trả lời ngắn không ghi đơn vị.
}
\end{ex}

% ===== Câu 18 =====
% ID: AIQ_L10C5_FULL_0342
% Mức: VD
\begin{ex}
% ID: AIQ_L10C5_FULL_0342
% Mức: VD
Trong bố trí thí nghiệm va chạm có 2 cổng quang điện độc lập. Có bao nhiêu thiết bị trực tiếp ghi thời điểm xe đi qua?
\shortans{2}
\loigiai{
Mỗi cổng quang điện là một thiết bị ghi thời điểm, nên có 2 thiết bị. Suy ra kết quả $2\ \mathrm{thiết\,bị}$. Đáp án trả lời ngắn không ghi đơn vị.
}
\end{ex}

% ===== Câu 19 =====
% ID: AIQ_L10C5_FULL_0343
% Mức: VD
\begin{ex}
% ID: AIQ_L10C5_FULL_0343
% Mức: VD
Trong bố trí thí nghiệm va chạm có 3 cổng quang điện độc lập. Có bao nhiêu thiết bị trực tiếp ghi thời điểm xe đi qua?
\shortans{3}
\loigiai{
Mỗi cổng quang điện là một thiết bị ghi thời điểm, nên có 3 thiết bị. Suy ra kết quả $3\ \mathrm{thiết\,bị}$. Đáp án trả lời ngắn không ghi đơn vị.
}
\end{ex}

% ===== Câu 20 =====
% ID: AIQ_L10C5_FULL_0344
% Mức: VDC
\begin{ex}
% ID: AIQ_L10C5_FULL_0344
% Mức: VDC
Trong bố trí thí nghiệm va chạm có 4 cổng quang điện độc lập. Có bao nhiêu thiết bị trực tiếp ghi thời điểm xe đi qua?
\shortans{4}
\loigiai{
Mỗi cổng quang điện là một thiết bị ghi thời điểm, nên có 4 thiết bị. Suy ra kết quả $4\ \mathrm{thiết\,bị}$. Đáp án trả lời ngắn không ghi đơn vị.
}
\end{ex}

% ===== Câu 21 =====
% ID: AIQ_L10C5_FULL_0345
% Mức: VDC
\begin{ex}
% ID: AIQ_L10C5_FULL_0345
% Mức: VDC
Trong bố trí thí nghiệm va chạm có 5 cổng quang điện độc lập. Có bao nhiêu thiết bị trực tiếp ghi thời điểm xe đi qua?
\shortans{5}
\loigiai{
Mỗi cổng quang điện là một thiết bị ghi thời điểm, nên có 5 thiết bị. Suy ra kết quả $5\ \mathrm{thiết\,bị}$. Đáp án trả lời ngắn không ghi đơn vị.
}
\end{ex}

% ===== Câu 22 =====
% ID: AIQ_L10C5_FULL_0346
% Mức: TH
\begin{bt}
% ID: AIQ_L10C5_FULL_0346
% Mức: TH
Trình bày cách giải: Trong bố trí thí nghiệm va chạm có 2 cổng quang điện độc lập. Có bao nhiêu thiết bị trực tiếp ghi thời điểm xe đi qua?
\loigiai{
Mỗi cổng quang điện là một thiết bị ghi thời điểm, nên có 2 thiết bị. Suy ra kết quả $2\ \mathrm{thiết\,bị}$.
}
\end{bt}

% ===== Câu 23 =====
% ID: AIQ_L10C5_FULL_0347
% Mức: TH
\begin{bt}
% ID: AIQ_L10C5_FULL_0347
% Mức: TH
Trình bày cách giải: Trong bố trí thí nghiệm va chạm có 3 cổng quang điện độc lập. Có bao nhiêu thiết bị trực tiếp ghi thời điểm xe đi qua?
\loigiai{
Mỗi cổng quang điện là một thiết bị ghi thời điểm, nên có 3 thiết bị. Suy ra kết quả $3\ \mathrm{thiết\,bị}$.
}
\end{bt}

% ===== Câu 24 =====
% ID: AIQ_L10C5_FULL_0348
% Mức: VD
\begin{bt}
% ID: AIQ_L10C5_FULL_0348
% Mức: VD
Trình bày cách giải: Trong bố trí thí nghiệm va chạm có 4 cổng quang điện độc lập. Có bao nhiêu thiết bị trực tiếp ghi thời điểm xe đi qua?
\loigiai{
Mỗi cổng quang điện là một thiết bị ghi thời điểm, nên có 4 thiết bị. Suy ra kết quả $4\ \mathrm{thiết\,bị}$.
}
\end{bt}

% ===== Câu 25 =====
% ID: AIQ_L10C5_FULL_0349
% Mức: VD
\begin{bt}
% ID: AIQ_L10C5_FULL_0349
% Mức: VD
Trình bày cách giải: Trong bố trí thí nghiệm va chạm có 5 cổng quang điện độc lập. Có bao nhiêu thiết bị trực tiếp ghi thời điểm xe đi qua?
\loigiai{
Mỗi cổng quang điện là một thiết bị ghi thời điểm, nên có 5 thiết bị. Suy ra kết quả $5\ \mathrm{thiết\,bị}$.
}
\end{bt}

% ===== Câu 26 =====
% ID: AIQ_L10C5_FULL_0350
% Mức: VDC
\begin{bt}
% ID: AIQ_L10C5_FULL_0350
% Mức: VDC
Trình bày cách giải: Trong bố trí thí nghiệm va chạm có 2 cổng quang điện độc lập. Có bao nhiêu thiết bị trực tiếp ghi thời điểm xe đi qua?
\loigiai{
Mỗi cổng quang điện là một thiết bị ghi thời điểm, nên có 2 thiết bị. Suy ra kết quả $2\ \mathrm{thiết\,bị}$.
}
\end{bt}

% ===== Câu 27 =====
% ID: AIQ_L10C5_FULL_0351
% Mức: VDC
\begin{bt}
% ID: AIQ_L10C5_FULL_0351
% Mức: VDC
Trình bày cách giải: Trong bố trí thí nghiệm va chạm có 3 cổng quang điện độc lập. Có bao nhiêu thiết bị trực tiếp ghi thời điểm xe đi qua?
\loigiai{
Mỗi cổng quang điện là một thiết bị ghi thời điểm, nên có 3 thiết bị. Suy ra kết quả $3\ \mathrm{thiết\,bị}$.
}
\end{bt}

% ===== Câu 28 =====
% ID: AIQ_L10C5_FULL_0352
% Mức: NB
\dangbt{Đo khối lượng, vận tốc và tính động lượng}
\begin{ex}
% ID: AIQ_L10C5_FULL_0352
% Mức: NB
Xe thí nghiệm có khối lượng $0,1$ kg, vận tốc đo được $0,4$ m/s. Tính động lượng.
\choice
{\True $0,04\ \mathrm{kg\,m/s}$}
{$0,5\ \mathrm{kg\,m/s}$}
{$0,4\ \mathrm{kg\,m/s}$}
{$0,3\ \mathrm{kg\,m/s}$}
\loigiai{
$p=mv=0,1\cdot0,4$. Suy ra kết quả $0,04\ \mathrm{kg\,m/s}$.
}
\end{ex}

% ===== Câu 29 =====
% ID: AIQ_L10C5_FULL_0353
% Mức: NB
\begin{ex}
% ID: AIQ_L10C5_FULL_0353
% Mức: NB
Xe thí nghiệm có khối lượng $0,15$ kg, vận tốc đo được $0,5$ m/s. Tính động lượng.
\choice
{$0,65\ \mathrm{kg\,m/s}$}
{\True $0,08\ \mathrm{kg\,m/s}$}
{$0,75\ \mathrm{kg\,m/s}$}
{$0,35\ \mathrm{kg\,m/s}$}
\loigiai{
$p=mv=0,15\cdot0,5$. Suy ra kết quả $0,08\ \mathrm{kg\,m/s}$.
}
\end{ex}

% ===== Câu 30 =====
% ID: AIQ_L10C5_FULL_0354
% Mức: TH
\begin{ex}
% ID: AIQ_L10C5_FULL_0354
% Mức: TH
Xe thí nghiệm có khối lượng $0,2$ kg, vận tốc đo được $0,6$ m/s. Tính động lượng.
\choice
{$0,8\ \mathrm{kg\,m/s}$}
{$1,2\ \mathrm{kg\,m/s}$}
{\True $0,12\ \mathrm{kg\,m/s}$}
{$0,4\ \mathrm{kg\,m/s}$}
\loigiai{
$p=mv=0,2\cdot0,6$. Suy ra kết quả $0,12\ \mathrm{kg\,m/s}$.
}
\end{ex}

% ===== Câu 31 =====
% ID: AIQ_L10C5_FULL_0355
% Mức: TH
\begin{ex}
% ID: AIQ_L10C5_FULL_0355
% Mức: TH
Xe thí nghiệm có khối lượng $0,25$ kg, vận tốc đo được $0,7$ m/s. Tính động lượng.
\choice
{$0,95\ \mathrm{kg\,m/s}$}
{$1,75\ \mathrm{kg\,m/s}$}
{$0,45\ \mathrm{kg\,m/s}$}
{\True $0,18\ \mathrm{kg\,m/s}$}
\loigiai{
$p=mv=0,25\cdot0,7$. Suy ra kết quả $0,18\ \mathrm{kg\,m/s}$.
}
\end{ex}

% ===== Câu 32 =====
% ID: AIQ_L10C5_FULL_0356
% Mức: VD
\begin{ex}
% ID: AIQ_L10C5_FULL_0356
% Mức: VD
Xe thí nghiệm có khối lượng $0,3$ kg, vận tốc đo được $0,8$ m/s. Tính động lượng.
\choice
{\True $0,24\ \mathrm{kg\,m/s}$}
{$1,1\ \mathrm{kg\,m/s}$}
{$2,4\ \mathrm{kg\,m/s}$}
{$0,5\ \mathrm{kg\,m/s}$}
\loigiai{
$p=mv=0,3\cdot0,8$. Suy ra kết quả $0,24\ \mathrm{kg\,m/s}$.
}
\end{ex}

% ===== Câu 33 =====
% ID: AIQ_L10C5_FULL_0357
% Mức: VD
\begin{ex}
% ID: AIQ_L10C5_FULL_0357
% Mức: VD
Xe thí nghiệm có khối lượng $0,1$ kg, vận tốc đo được $0,9$ m/s. Tính động lượng.
\choice
{$1\ \mathrm{kg\,m/s}$}
{\True $0,09\ \mathrm{kg\,m/s}$}
{$0,9\ \mathrm{kg\,m/s}$}
{$0,8\ \mathrm{kg\,m/s}$}
\loigiai{
$p=mv=0,1\cdot0,9$. Suy ra kết quả $0,09\ \mathrm{kg\,m/s}$.
}
\end{ex}

% ===== Câu 34 =====
% ID: AIQ_L10C5_FULL_0358
% Mức: VD
\begin{ex}
% ID: AIQ_L10C5_FULL_0358
% Mức: VD
Xe thí nghiệm có khối lượng $0,15$ kg, vận tốc đo được $1$ m/s. Tính động lượng.
\choice
{$1,15\ \mathrm{kg\,m/s}$}
{$1,5\ \mathrm{kg\,m/s}$}
{\True $0,15\ \mathrm{kg\,m/s}$}
{$0,85\ \mathrm{kg\,m/s}$}
\loigiai{
$p=mv=0,15\cdot1$. Suy ra kết quả $0,15\ \mathrm{kg\,m/s}$.
}
\end{ex}

% ===== Câu 35 =====
% ID: AIQ_L10C5_FULL_0359
% Mức: VDC
\begin{ex}
% ID: AIQ_L10C5_FULL_0359
% Mức: VDC
Xe thí nghiệm có khối lượng $0,2$ kg, vận tốc đo được $1,1$ m/s. Tính động lượng.
\choice
{$1,3\ \mathrm{kg\,m/s}$}
{$2,2\ \mathrm{kg\,m/s}$}
{$0,9\ \mathrm{kg\,m/s}$}
{\True $0,22\ \mathrm{kg\,m/s}$}
\loigiai{
$p=mv=0,2\cdot1,1$. Suy ra kết quả $0,22\ \mathrm{kg\,m/s}$.
}
\end{ex}

% ===== Câu 36 =====
% ID: AIQ_L10C5_FULL_0360
% Mức: VDC
\begin{ex}
% ID: AIQ_L10C5_FULL_0360
% Mức: VDC
Xe thí nghiệm có khối lượng $0,25$ kg, vận tốc đo được $1,2$ m/s. Tính động lượng.
\choice
{\True $0,3\ \mathrm{kg\,m/s}$}
{$1,45\ \mathrm{kg\,m/s}$}
{$3\ \mathrm{kg\,m/s}$}
{$0,95\ \mathrm{kg\,m/s}$}
\loigiai{
$p=mv=0,25\cdot1,2$. Suy ra kết quả $0,3\ \mathrm{kg\,m/s}$.
}
\end{ex}

% ===== Câu 37 =====
% ID: AIQ_L10C5_FULL_0361
% Mức: VDC
\begin{ex}
% ID: AIQ_L10C5_FULL_0361
% Mức: VDC
Xe thí nghiệm có khối lượng $0,3$ kg, vận tốc đo được $1,3$ m/s. Tính động lượng.
\choice
{$1,6\ \mathrm{kg\,m/s}$}
{\True $0,39\ \mathrm{kg\,m/s}$}
{$3,9\ \mathrm{kg\,m/s}$}
{$1\ \mathrm{kg\,m/s}$}
\loigiai{
$p=mv=0,3\cdot1,3$. Suy ra kết quả $0,39\ \mathrm{kg\,m/s}$.
}
\end{ex}

% ===== Câu 38 =====
% ID: AIQ_L10C5_FULL_0362
% Mức: NB
\begin{ex}
% ID: AIQ_L10C5_FULL_0362
% Mức: NB
Xét các phát biểu thuộc dạng “Đo khối lượng, vận tốc và tính động lượng” ở tình huống số 1.
\choiceTF
{\True Động lượng tính bằng $p=mv$.}
{\True Khối lượng phải đổi về kg.}
{Vận tốc không cần xác định chiều.}
{\True Số đo cần ghi đơn vị.}
\loigiai{
\begin{itemchoice}
\item (Đúng) Động lượng tính bằng $p=mv$. (Vì): Phù hợp với định nghĩa hoặc hệ thức vật lí. b) Đúng: Phù hợp với định nghĩa hoặc hệ thức vật lí. c) Sai: Không phù hợp với định nghĩa hoặc hệ thức vật lí. d) Đúng: Phù hợp với định nghĩa hoặc hệ thức vật lí
\item (Đúng) Khối lượng phải đổi về kg.
\item (Sai) Vận tốc không cần xác định chiều.
\item (Đúng) Số đo cần ghi đơn vị.
\end{itemchoice}
}
\end{ex}

% ===== Câu 39 =====
% ID: AIQ_L10C5_FULL_0363
% Mức: TH
\begin{ex}
% ID: AIQ_L10C5_FULL_0363
% Mức: TH
Xét các phát biểu thuộc dạng “Đo khối lượng, vận tốc và tính động lượng” ở tình huống số 2.
\choiceTF
{\True Động lượng tính bằng $p=mv$.}
{\True Khối lượng phải đổi về kg.}
{Vận tốc không cần xác định chiều.}
{\True Số đo cần ghi đơn vị.}
\loigiai{
\begin{itemchoice}
\item (Đúng) Động lượng tính bằng $p=mv$. (Vì): Phù hợp với định nghĩa hoặc hệ thức vật lí. b) Đúng: Phù hợp với định nghĩa hoặc hệ thức vật lí. c) Sai: Không phù hợp với định nghĩa hoặc hệ thức vật lí. d) Đúng: Phù hợp với định nghĩa hoặc hệ thức vật lí
\item (Đúng) Khối lượng phải đổi về kg.
\item (Sai) Vận tốc không cần xác định chiều.
\item (Đúng) Số đo cần ghi đơn vị.
\end{itemchoice}
}
\end{ex}

% ===== Câu 40 =====
% ID: AIQ_L10C5_FULL_0364
% Mức: TH
\begin{ex}
% ID: AIQ_L10C5_FULL_0364
% Mức: TH
Xét các phát biểu thuộc dạng “Đo khối lượng, vận tốc và tính động lượng” ở tình huống số 3.
\choiceTF
{\True Động lượng tính bằng $p=mv$.}
{\True Khối lượng phải đổi về kg.}
{Vận tốc không cần xác định chiều.}
{\True Số đo cần ghi đơn vị.}
\loigiai{
\begin{itemchoice}
\item (Đúng) Động lượng tính bằng $p=mv$. (Vì): Phù hợp với định nghĩa hoặc hệ thức vật lí. b) Đúng: Phù hợp với định nghĩa hoặc hệ thức vật lí. c) Sai: Không phù hợp với định nghĩa hoặc hệ thức vật lí. d) Đúng: Phù hợp với định nghĩa hoặc hệ thức vật lí
\item (Đúng) Khối lượng phải đổi về kg.
\item (Sai) Vận tốc không cần xác định chiều.
\item (Đúng) Số đo cần ghi đơn vị.
\end{itemchoice}
}
\end{ex}

% ===== Câu 41 =====
% ID: AIQ_L10C5_FULL_0365
% Mức: VD
\begin{ex}
% ID: AIQ_L10C5_FULL_0365
% Mức: VD
Xét các phát biểu thuộc dạng “Đo khối lượng, vận tốc và tính động lượng” ở tình huống số 4.
\choiceTF
{\True Động lượng tính bằng $p=mv$.}
{\True Khối lượng phải đổi về kg.}
{Vận tốc không cần xác định chiều.}
{\True Số đo cần ghi đơn vị.}
\loigiai{
\begin{itemchoice}
\item (Đúng) Động lượng tính bằng $p=mv$. (Vì): Phù hợp với định nghĩa hoặc hệ thức vật lí. b) Đúng: Phù hợp với định nghĩa hoặc hệ thức vật lí. c) Sai: Không phù hợp với định nghĩa hoặc hệ thức vật lí. d) Đúng: Phù hợp với định nghĩa hoặc hệ thức vật lí
\item (Đúng) Khối lượng phải đổi về kg.
\item (Sai) Vận tốc không cần xác định chiều.
\item (Đúng) Số đo cần ghi đơn vị.
\end{itemchoice}
}
\end{ex}

% ===== Câu 42 =====
% ID: AIQ_L10C5_FULL_0366
% Mức: VDC
\begin{ex}
% ID: AIQ_L10C5_FULL_0366
% Mức: VDC
Xét các phát biểu thuộc dạng “Đo khối lượng, vận tốc và tính động lượng” ở tình huống số 5.
\choiceTF
{\True Động lượng tính bằng $p=mv$.}
{\True Khối lượng phải đổi về kg.}
{Vận tốc không cần xác định chiều.}
{\True Số đo cần ghi đơn vị.}
\loigiai{
\begin{itemchoice}
\item (Đúng) Động lượng tính bằng $p=mv$. (Vì): Phù hợp với định nghĩa hoặc hệ thức vật lí. b) Đúng: Phù hợp với định nghĩa hoặc hệ thức vật lí. c) Sai: Không phù hợp với định nghĩa hoặc hệ thức vật lí. d) Đúng: Phù hợp với định nghĩa hoặc hệ thức vật lí
\item (Đúng) Khối lượng phải đổi về kg.
\item (Sai) Vận tốc không cần xác định chiều.
\item (Đúng) Số đo cần ghi đơn vị.
\end{itemchoice}
}
\end{ex}

% ===== Câu 43 =====
% ID: AIQ_L10C5_FULL_0367
% Mức: TH
\begin{ex}
% ID: AIQ_L10C5_FULL_0367
% Mức: TH
Xe thí nghiệm có khối lượng $0,1$ kg, vận tốc đo được $1,4$ m/s. Tính động lượng.
\shortans{0,14}
\loigiai{
$p=mv=0,1\cdot1,4$. Suy ra kết quả $0,14\ \mathrm{kg\,m/s}$. Đáp án trả lời ngắn không ghi đơn vị.
}
\end{ex}

% ===== Câu 44 =====
% ID: AIQ_L10C5_FULL_0368
% Mức: TH
\begin{ex}
% ID: AIQ_L10C5_FULL_0368
% Mức: TH
Xe thí nghiệm có khối lượng $0,15$ kg, vận tốc đo được $1,5$ m/s. Tính động lượng.
\shortans{0,23}
\loigiai{
$p=mv=0,15\cdot1,5$. Suy ra kết quả $0,23\ \mathrm{kg\,m/s}$. Đáp án trả lời ngắn không ghi đơn vị.
}
\end{ex}

% ===== Câu 45 =====
% ID: AIQ_L10C5_FULL_0369
% Mức: VD
\begin{ex}
% ID: AIQ_L10C5_FULL_0369
% Mức: VD
Xe thí nghiệm có khối lượng $0,2$ kg, vận tốc đo được $1,6$ m/s. Tính động lượng.
\shortans{0,32}
\loigiai{
$p=mv=0,2\cdot1,6$. Suy ra kết quả $0,32\ \mathrm{kg\,m/s}$. Đáp án trả lời ngắn không ghi đơn vị.
}
\end{ex}

% ===== Câu 46 =====
% ID: AIQ_L10C5_FULL_0370
% Mức: VD
\begin{ex}
% ID: AIQ_L10C5_FULL_0370
% Mức: VD
Xe thí nghiệm có khối lượng $0,25$ kg, vận tốc đo được $1,7$ m/s. Tính động lượng.
\shortans{0,43}
\loigiai{
$p=mv=0,25\cdot1,7$. Suy ra kết quả $0,43\ \mathrm{kg\,m/s}$. Đáp án trả lời ngắn không ghi đơn vị.
}
\end{ex}

% ===== Câu 47 =====
% ID: AIQ_L10C5_FULL_0371
% Mức: VDC
\begin{ex}
% ID: AIQ_L10C5_FULL_0371
% Mức: VDC
Xe thí nghiệm có khối lượng $0,3$ kg, vận tốc đo được $1,8$ m/s. Tính động lượng.
\shortans{0,54}
\loigiai{
$p=mv=0,3\cdot1,8$. Suy ra kết quả $0,54\ \mathrm{kg\,m/s}$. Đáp án trả lời ngắn không ghi đơn vị.
}
\end{ex}

% ===== Câu 48 =====
% ID: AIQ_L10C5_FULL_0372
% Mức: VDC
\begin{ex}
% ID: AIQ_L10C5_FULL_0372
% Mức: VDC
Xe thí nghiệm có khối lượng $0,1$ kg, vận tốc đo được $1,9$ m/s. Tính động lượng.
\shortans{0,19}
\loigiai{
$p=mv=0,1\cdot1,9$. Suy ra kết quả $0,19\ \mathrm{kg\,m/s}$. Đáp án trả lời ngắn không ghi đơn vị.
}
\end{ex}

% ===== Câu 49 =====
% ID: AIQ_L10C5_FULL_0373
% Mức: TH
\begin{bt}
% ID: AIQ_L10C5_FULL_0373
% Mức: TH
Trình bày cách giải: Xe thí nghiệm có khối lượng $0,15$ kg, vận tốc đo được $2$ m/s. Tính động lượng.
\loigiai{
$p=mv=0,15\cdot2$. Suy ra kết quả $0,3\ \mathrm{kg\,m/s}$.
}
\end{bt}

% ===== Câu 50 =====
% ID: AIQ_L10C5_FULL_0374
% Mức: TH
\begin{bt}
% ID: AIQ_L10C5_FULL_0374
% Mức: TH
Trình bày cách giải: Xe thí nghiệm có khối lượng $0,2$ kg, vận tốc đo được $2,1$ m/s. Tính động lượng.
\loigiai{
$p=mv=0,2\cdot2,1$. Suy ra kết quả $0,42\ \mathrm{kg\,m/s}$.
}
\end{bt}

% ===== Câu 51 =====
% ID: AIQ_L10C5_FULL_0375
% Mức: VD
\begin{bt}
% ID: AIQ_L10C5_FULL_0375
% Mức: VD
Trình bày cách giải: Xe thí nghiệm có khối lượng $0,25$ kg, vận tốc đo được $2,2$ m/s. Tính động lượng.
\loigiai{
$p=mv=0,25\cdot2,2$. Suy ra kết quả $0,55\ \mathrm{kg\,m/s}$.
}
\end{bt}

% ===== Câu 52 =====
% ID: AIQ_L10C5_FULL_0376
% Mức: VD
\begin{bt}
% ID: AIQ_L10C5_FULL_0376
% Mức: VD
Trình bày cách giải: Xe thí nghiệm có khối lượng $0,3$ kg, vận tốc đo được $2,3$ m/s. Tính động lượng.
\loigiai{
$p=mv=0,3\cdot2,3$. Suy ra kết quả $0,69\ \mathrm{kg\,m/s}$.
}
\end{bt}

% ===== Câu 53 =====
% ID: AIQ_L10C5_FULL_0377
% Mức: VDC
\begin{bt}
% ID: AIQ_L10C5_FULL_0377
% Mức: VDC
Trình bày cách giải: Xe thí nghiệm có khối lượng $0,1$ kg, vận tốc đo được $2,4$ m/s. Tính động lượng.
\loigiai{
$p=mv=0,1\cdot2,4$. Suy ra kết quả $0,24\ \mathrm{kg\,m/s}$.
}
\end{bt}

% ===== Câu 54 =====
% ID: AIQ_L10C5_FULL_0378
% Mức: VDC
\begin{bt}
% ID: AIQ_L10C5_FULL_0378
% Mức: VDC
Trình bày cách giải: Xe thí nghiệm có khối lượng $0,15$ kg, vận tốc đo được $2,5$ m/s. Tính động lượng.
\loigiai{
$p=mv=0,15\cdot2,5$. Suy ra kết quả $0,38\ \mathrm{kg\,m/s}$.
}
\end{bt}

% ===== Câu 55 =====
% ID: AIQ_L10C5_FULL_0379
% Mức: NB
\dangbt{Xử lí số liệu động lượng trước và sau va chạm}
\begin{ex}
% ID: AIQ_L10C5_FULL_0379
% Mức: NB
Trước va chạm, hai xe chuyển động cùng chiều với động lượng $0,2$ và $0,1\ \mathrm{kg\,m/s}$. Tính tổng động lượng.
\choice
{\True $0,3\ \mathrm{kg\,m/s}$}
{$0,1\ \mathrm{kg\,m/s}$}
{$0,02\ \mathrm{kg\,m/s}$}
{$0,6\ \mathrm{kg\,m/s}$}
\loigiai{
Cùng chiều nên $p=p_1+p_2=0,2+0,1$. Suy ra kết quả $0,3\ \mathrm{kg\,m/s}$.
}
\end{ex}

% ===== Câu 56 =====
% ID: AIQ_L10C5_FULL_0380
% Mức: NB
\begin{ex}
% ID: AIQ_L10C5_FULL_0380
% Mức: NB
Trước va chạm, hai xe chuyển động cùng chiều với động lượng $0,22$ và $0,11\ \mathrm{kg\,m/s}$. Tính tổng động lượng.
\choice
{$0,11\ \mathrm{kg\,m/s}$}
{\True $0,33\ \mathrm{kg\,m/s}$}
{$0,02\ \mathrm{kg\,m/s}$}
{$0,66\ \mathrm{kg\,m/s}$}
\loigiai{
Cùng chiều nên $p=p_1+p_2=0,22+0,11$. Suy ra kết quả $0,33\ \mathrm{kg\,m/s}$.
}
\end{ex}

% ===== Câu 57 =====
% ID: AIQ_L10C5_FULL_0381
% Mức: TH
\begin{ex}
% ID: AIQ_L10C5_FULL_0381
% Mức: TH
Trước va chạm, hai xe chuyển động cùng chiều với động lượng $0,24$ và $0,12\ \mathrm{kg\,m/s}$. Tính tổng động lượng.
\choice
{$0,12\ \mathrm{kg\,m/s}$}
{$0,03\ \mathrm{kg\,m/s}$}
{\True $0,36\ \mathrm{kg\,m/s}$}
{$0,72\ \mathrm{kg\,m/s}$}
\loigiai{
Cùng chiều nên $p=p_1+p_2=0,24+0,12$. Suy ra kết quả $0,36\ \mathrm{kg\,m/s}$.
}
\end{ex}

% ===== Câu 58 =====
% ID: AIQ_L10C5_FULL_0382
% Mức: TH
\begin{ex}
% ID: AIQ_L10C5_FULL_0382
% Mức: TH
Trước va chạm, hai xe chuyển động cùng chiều với động lượng $0,26$ và $0,13\ \mathrm{kg\,m/s}$. Tính tổng động lượng.
\choice
{$0,13\ \mathrm{kg\,m/s}$}
{$0,03\ \mathrm{kg\,m/s}$}
{$0,78\ \mathrm{kg\,m/s}$}
{\True $0,39\ \mathrm{kg\,m/s}$}
\loigiai{
Cùng chiều nên $p=p_1+p_2=0,26+0,13$. Suy ra kết quả $0,39\ \mathrm{kg\,m/s}$.
}
\end{ex}

% ===== Câu 59 =====
% ID: AIQ_L10C5_FULL_0383
% Mức: VD
\begin{ex}
% ID: AIQ_L10C5_FULL_0383
% Mức: VD
Trước va chạm, hai xe chuyển động cùng chiều với động lượng $0,28$ và $0,1\ \mathrm{kg\,m/s}$. Tính tổng động lượng.
\choice
{\True $0,38\ \mathrm{kg\,m/s}$}
{$0,18\ \mathrm{kg\,m/s}$}
{$0,03\ \mathrm{kg\,m/s}$}
{$0,76\ \mathrm{kg\,m/s}$}
\loigiai{
Cùng chiều nên $p=p_1+p_2=0,28+0,1$. Suy ra kết quả $0,38\ \mathrm{kg\,m/s}$.
}
\end{ex}

% ===== Câu 60 =====
% ID: AIQ_L10C5_FULL_0384
% Mức: VD
\begin{ex}
% ID: AIQ_L10C5_FULL_0384
% Mức: VD
Trước va chạm, hai xe chuyển động cùng chiều với động lượng $0,3$ và $0,11\ \mathrm{kg\,m/s}$. Tính tổng động lượng.
\choice
{$0,19\ \mathrm{kg\,m/s}$}
{\True $0,41\ \mathrm{kg\,m/s}$}
{$0,03\ \mathrm{kg\,m/s}$}
{$0,82\ \mathrm{kg\,m/s}$}
\loigiai{
Cùng chiều nên $p=p_1+p_2=0,3+0,11$. Suy ra kết quả $0,41\ \mathrm{kg\,m/s}$.
}
\end{ex}

% ===== Câu 61 =====
% ID: AIQ_L10C5_FULL_0385
% Mức: VD
\begin{ex}
% ID: AIQ_L10C5_FULL_0385
% Mức: VD
Trước va chạm, hai xe chuyển động cùng chiều với động lượng $0,32$ và $0,12\ \mathrm{kg\,m/s}$. Tính tổng động lượng.
\choice
{$0,2\ \mathrm{kg\,m/s}$}
{$0,04\ \mathrm{kg\,m/s}$}
{\True $0,44\ \mathrm{kg\,m/s}$}
{$0,88\ \mathrm{kg\,m/s}$}
\loigiai{
Cùng chiều nên $p=p_1+p_2=0,32+0,12$. Suy ra kết quả $0,44\ \mathrm{kg\,m/s}$.
}
\end{ex}

% ===== Câu 62 =====
% ID: AIQ_L10C5_FULL_0386
% Mức: VDC
\begin{ex}
% ID: AIQ_L10C5_FULL_0386
% Mức: VDC
Trước va chạm, hai xe chuyển động cùng chiều với động lượng $0,34$ và $0,13\ \mathrm{kg\,m/s}$. Tính tổng động lượng.
\choice
{$0,21\ \mathrm{kg\,m/s}$}
{$0,04\ \mathrm{kg\,m/s}$}
{$0,94\ \mathrm{kg\,m/s}$}
{\True $0,47\ \mathrm{kg\,m/s}$}
\loigiai{
Cùng chiều nên $p=p_1+p_2=0,34+0,13$. Suy ra kết quả $0,47\ \mathrm{kg\,m/s}$.
}
\end{ex}

% ===== Câu 63 =====
% ID: AIQ_L10C5_FULL_0387
% Mức: VDC
\begin{ex}
% ID: AIQ_L10C5_FULL_0387
% Mức: VDC
Trước va chạm, hai xe chuyển động cùng chiều với động lượng $0,36$ và $0,1\ \mathrm{kg\,m/s}$. Tính tổng động lượng.
\choice
{\True $0,46\ \mathrm{kg\,m/s}$}
{$0,26\ \mathrm{kg\,m/s}$}
{$0,04\ \mathrm{kg\,m/s}$}
{$0,92\ \mathrm{kg\,m/s}$}
\loigiai{
Cùng chiều nên $p=p_1+p_2=0,36+0,1$. Suy ra kết quả $0,46\ \mathrm{kg\,m/s}$.
}
\end{ex}

% ===== Câu 64 =====
% ID: AIQ_L10C5_FULL_0388
% Mức: VDC
\begin{ex}
% ID: AIQ_L10C5_FULL_0388
% Mức: VDC
Trước va chạm, hai xe chuyển động cùng chiều với động lượng $0,38$ và $0,11\ \mathrm{kg\,m/s}$. Tính tổng động lượng.
\choice
{$0,27\ \mathrm{kg\,m/s}$}
{\True $0,49\ \mathrm{kg\,m/s}$}
{$0,04\ \mathrm{kg\,m/s}$}
{$0,98\ \mathrm{kg\,m/s}$}
\loigiai{
Cùng chiều nên $p=p_1+p_2=0,38+0,11$. Suy ra kết quả $0,49\ \mathrm{kg\,m/s}$.
}
\end{ex}

% ===== Câu 65 =====
% ID: AIQ_L10C5_FULL_0389
% Mức: NB
\begin{ex}
% ID: AIQ_L10C5_FULL_0389
% Mức: NB
Xét các phát biểu thuộc dạng “Xử lí số liệu động lượng trước và sau va chạm” ở tình huống số 1.
\choiceTF
{\True Phải chọn chiều dương khi xử lí số liệu.}
{\True Động lượng ngược chiều mang dấu trái nhau.}
{Chỉ cần so sánh độ lớn từng xe.}
{\True Tổng động lượng được tính bằng tổng đại số có dấu trong chuyển động thẳng.}
\loigiai{
\begin{itemchoice}
\item (Đúng) Phải chọn chiều dương khi xử lí số liệu. (Vì): Phù hợp với định nghĩa hoặc hệ thức vật lí. b) Đúng: Phù hợp với định nghĩa hoặc hệ thức vật lí. c) Sai: Không phù hợp với định nghĩa hoặc hệ thức vật lí. d) Đúng: Phù hợp với định nghĩa hoặc hệ thức vật lí
\item (Đúng) Động lượng ngược chiều mang dấu trái nhau.
\item (Sai) Chỉ cần so sánh độ lớn từng xe.
\item (Đúng) Tổng động lượng được tính bằng tổng đại số có dấu trong chuyển động thẳng.
\end{itemchoice}
}
\end{ex}

% ===== Câu 66 =====
% ID: AIQ_L10C5_FULL_0390
% Mức: TH
\begin{ex}
% ID: AIQ_L10C5_FULL_0390
% Mức: TH
Xét các phát biểu thuộc dạng “Xử lí số liệu động lượng trước và sau va chạm” ở tình huống số 2.
\choiceTF
{\True Phải chọn chiều dương khi xử lí số liệu.}
{\True Động lượng ngược chiều mang dấu trái nhau.}
{Chỉ cần so sánh độ lớn từng xe.}
{\True Tổng động lượng được tính bằng tổng đại số có dấu trong chuyển động thẳng.}
\loigiai{
\begin{itemchoice}
\item (Đúng) Phải chọn chiều dương khi xử lí số liệu. (Vì): Phù hợp với định nghĩa hoặc hệ thức vật lí. b) Đúng: Phù hợp với định nghĩa hoặc hệ thức vật lí. c) Sai: Không phù hợp với định nghĩa hoặc hệ thức vật lí. d) Đúng: Phù hợp với định nghĩa hoặc hệ thức vật lí
\item (Đúng) Động lượng ngược chiều mang dấu trái nhau.
\item (Sai) Chỉ cần so sánh độ lớn từng xe.
\item (Đúng) Tổng động lượng được tính bằng tổng đại số có dấu trong chuyển động thẳng.
\end{itemchoice}
}
\end{ex}

% ===== Câu 67 =====
% ID: AIQ_L10C5_FULL_0391
% Mức: TH
\begin{ex}
% ID: AIQ_L10C5_FULL_0391
% Mức: TH
Xét các phát biểu thuộc dạng “Xử lí số liệu động lượng trước và sau va chạm” ở tình huống số 3.
\choiceTF
{\True Phải chọn chiều dương khi xử lí số liệu.}
{\True Động lượng ngược chiều mang dấu trái nhau.}
{Chỉ cần so sánh độ lớn từng xe.}
{\True Tổng động lượng được tính bằng tổng đại số có dấu trong chuyển động thẳng.}
\loigiai{
\begin{itemchoice}
\item (Đúng) Phải chọn chiều dương khi xử lí số liệu. (Vì): Phù hợp với định nghĩa hoặc hệ thức vật lí. b) Đúng: Phù hợp với định nghĩa hoặc hệ thức vật lí. c) Sai: Không phù hợp với định nghĩa hoặc hệ thức vật lí. d) Đúng: Phù hợp với định nghĩa hoặc hệ thức vật lí
\item (Đúng) Động lượng ngược chiều mang dấu trái nhau.
\item (Sai) Chỉ cần so sánh độ lớn từng xe.
\item (Đúng) Tổng động lượng được tính bằng tổng đại số có dấu trong chuyển động thẳng.
\end{itemchoice}
}
\end{ex}

% ===== Câu 68 =====
% ID: AIQ_L10C5_FULL_0392
% Mức: VD
\begin{ex}
% ID: AIQ_L10C5_FULL_0392
% Mức: VD
Xét các phát biểu thuộc dạng “Xử lí số liệu động lượng trước và sau va chạm” ở tình huống số 4.
\choiceTF
{\True Phải chọn chiều dương khi xử lí số liệu.}
{\True Động lượng ngược chiều mang dấu trái nhau.}
{Chỉ cần so sánh độ lớn từng xe.}
{\True Tổng động lượng được tính bằng tổng đại số có dấu trong chuyển động thẳng.}
\loigiai{
\begin{itemchoice}
\item (Đúng) Phải chọn chiều dương khi xử lí số liệu. (Vì): Phù hợp với định nghĩa hoặc hệ thức vật lí. b) Đúng: Phù hợp với định nghĩa hoặc hệ thức vật lí. c) Sai: Không phù hợp với định nghĩa hoặc hệ thức vật lí. d) Đúng: Phù hợp với định nghĩa hoặc hệ thức vật lí
\item (Đúng) Động lượng ngược chiều mang dấu trái nhau.
\item (Sai) Chỉ cần so sánh độ lớn từng xe.
\item (Đúng) Tổng động lượng được tính bằng tổng đại số có dấu trong chuyển động thẳng.
\end{itemchoice}
}
\end{ex}

% ===== Câu 69 =====
% ID: AIQ_L10C5_FULL_0393
% Mức: VDC
\begin{ex}
% ID: AIQ_L10C5_FULL_0393
% Mức: VDC
Xét các phát biểu thuộc dạng “Xử lí số liệu động lượng trước và sau va chạm” ở tình huống số 5.
\choiceTF
{\True Phải chọn chiều dương khi xử lí số liệu.}
{\True Động lượng ngược chiều mang dấu trái nhau.}
{Chỉ cần so sánh độ lớn từng xe.}
{\True Tổng động lượng được tính bằng tổng đại số có dấu trong chuyển động thẳng.}
\loigiai{
\begin{itemchoice}
\item (Đúng) Phải chọn chiều dương khi xử lí số liệu. (Vì): Phù hợp với định nghĩa hoặc hệ thức vật lí. b) Đúng: Phù hợp với định nghĩa hoặc hệ thức vật lí. c) Sai: Không phù hợp với định nghĩa hoặc hệ thức vật lí. d) Đúng: Phù hợp với định nghĩa hoặc hệ thức vật lí
\item (Đúng) Động lượng ngược chiều mang dấu trái nhau.
\item (Sai) Chỉ cần so sánh độ lớn từng xe.
\item (Đúng) Tổng động lượng được tính bằng tổng đại số có dấu trong chuyển động thẳng.
\end{itemchoice}
}
\end{ex}

% ===== Câu 70 =====
% ID: AIQ_L10C5_FULL_0394
% Mức: TH
\begin{ex}
% ID: AIQ_L10C5_FULL_0394
% Mức: TH
Trước va chạm, hai xe chuyển động cùng chiều với động lượng $0,4$ và $0,12\ \mathrm{kg\,m/s}$. Tính tổng động lượng.
\shortans{0,52}
\loigiai{
Cùng chiều nên $p=p_1+p_2=0,4+0,12$. Suy ra kết quả $0,52\ \mathrm{kg\,m/s}$. Đáp án trả lời ngắn không ghi đơn vị.
}
\end{ex}

% ===== Câu 71 =====
% ID: AIQ_L10C5_FULL_0395
% Mức: TH
\begin{ex}
% ID: AIQ_L10C5_FULL_0395
% Mức: TH
Trước va chạm, hai xe chuyển động cùng chiều với động lượng $0,42$ và $0,13\ \mathrm{kg\,m/s}$. Tính tổng động lượng.
\shortans{0,55}
\loigiai{
Cùng chiều nên $p=p_1+p_2=0,42+0,13$. Suy ra kết quả $0,55\ \mathrm{kg\,m/s}$. Đáp án trả lời ngắn không ghi đơn vị.
}
\end{ex}

% ===== Câu 72 =====
% ID: AIQ_L10C5_FULL_0396
% Mức: VD
\begin{ex}
% ID: AIQ_L10C5_FULL_0396
% Mức: VD
Trước va chạm, hai xe chuyển động cùng chiều với động lượng $0,44$ và $0,1\ \mathrm{kg\,m/s}$. Tính tổng động lượng.
\shortans{0,54}
\loigiai{
Cùng chiều nên $p=p_1+p_2=0,44+0,1$. Suy ra kết quả $0,54\ \mathrm{kg\,m/s}$. Đáp án trả lời ngắn không ghi đơn vị.
}
\end{ex}

% ===== Câu 73 =====
% ID: AIQ_L10C5_FULL_0397
% Mức: VD
\begin{ex}
% ID: AIQ_L10C5_FULL_0397
% Mức: VD
Trước va chạm, hai xe chuyển động cùng chiều với động lượng $0,46$ và $0,11\ \mathrm{kg\,m/s}$. Tính tổng động lượng.
\shortans{0,57}
\loigiai{
Cùng chiều nên $p=p_1+p_2=0,46+0,11$. Suy ra kết quả $0,57\ \mathrm{kg\,m/s}$. Đáp án trả lời ngắn không ghi đơn vị.
}
\end{ex}

% ===== Câu 74 =====
% ID: AIQ_L10C5_FULL_0398
% Mức: VDC
\begin{ex}
% ID: AIQ_L10C5_FULL_0398
% Mức: VDC
Trước va chạm, hai xe chuyển động cùng chiều với động lượng $0,48$ và $0,12\ \mathrm{kg\,m/s}$. Tính tổng động lượng.
\shortans{0,6}
\loigiai{
Cùng chiều nên $p=p_1+p_2=0,48+0,12$. Suy ra kết quả $0,6\ \mathrm{kg\,m/s}$. Đáp án trả lời ngắn không ghi đơn vị.
}
\end{ex}

% ===== Câu 75 =====
% ID: AIQ_L10C5_FULL_0399
% Mức: VDC
\begin{ex}
% ID: AIQ_L10C5_FULL_0399
% Mức: VDC
Trước va chạm, hai xe chuyển động cùng chiều với động lượng $0,5$ và $0,13\ \mathrm{kg\,m/s}$. Tính tổng động lượng.
\shortans{0,63}
\loigiai{
Cùng chiều nên $p=p_1+p_2=0,5+0,13$. Suy ra kết quả $0,63\ \mathrm{kg\,m/s}$. Đáp án trả lời ngắn không ghi đơn vị.
}
\end{ex}

% ===== Câu 76 =====
% ID: AIQ_L10C5_FULL_0400
% Mức: TH
\begin{bt}
% ID: AIQ_L10C5_FULL_0400
% Mức: TH
Trình bày cách giải: Trước va chạm, hai xe chuyển động cùng chiều với động lượng $0,52$ và $0,1\ \mathrm{kg\,m/s}$. Tính tổng động lượng.
\loigiai{
Cùng chiều nên $p=p_1+p_2=0,52+0,1$. Suy ra kết quả $0,62\ \mathrm{kg\,m/s}$.
}
\end{bt}

% ===== Câu 77 =====
% ID: AIQ_L10C5_FULL_0401
% Mức: TH
\begin{bt}
% ID: AIQ_L10C5_FULL_0401
% Mức: TH
Trình bày cách giải: Trước va chạm, hai xe chuyển động cùng chiều với động lượng $0,54$ và $0,11\ \mathrm{kg\,m/s}$. Tính tổng động lượng.
\loigiai{
Cùng chiều nên $p=p_1+p_2=0,54+0,11$. Suy ra kết quả $0,65\ \mathrm{kg\,m/s}$.
}
\end{bt}

% ===== Câu 78 =====
% ID: AIQ_L10C5_FULL_0402
% Mức: VD
\begin{bt}
% ID: AIQ_L10C5_FULL_0402
% Mức: VD
Trình bày cách giải: Trước va chạm, hai xe chuyển động cùng chiều với động lượng $0,56$ và $0,12\ \mathrm{kg\,m/s}$. Tính tổng động lượng.
\loigiai{
Cùng chiều nên $p=p_1+p_2=0,56+0,12$. Suy ra kết quả $0,68\ \mathrm{kg\,m/s}$.
}
\end{bt}

% ===== Câu 79 =====
% ID: AIQ_L10C5_FULL_0403
% Mức: VD
\begin{bt}
% ID: AIQ_L10C5_FULL_0403
% Mức: VD
Trình bày cách giải: Trước va chạm, hai xe chuyển động cùng chiều với động lượng $0,58$ và $0,13\ \mathrm{kg\,m/s}$. Tính tổng động lượng.
\loigiai{
Cùng chiều nên $p=p_1+p_2=0,58+0,13$. Suy ra kết quả $0,71\ \mathrm{kg\,m/s}$.
}
\end{bt}

% ===== Câu 80 =====
% ID: AIQ_L10C5_FULL_0404
% Mức: VDC
\begin{bt}
% ID: AIQ_L10C5_FULL_0404
% Mức: VDC
Trình bày cách giải: Trước va chạm, hai xe chuyển động cùng chiều với động lượng $0,6$ và $0,1\ \mathrm{kg\,m/s}$. Tính tổng động lượng.
\loigiai{
Cùng chiều nên $p=p_1+p_2=0,6+0,1$. Suy ra kết quả $0,7\ \mathrm{kg\,m/s}$.
}
\end{bt}

% ===== Câu 81 =====
% ID: AIQ_L10C5_FULL_0405
% Mức: VDC
\begin{bt}
% ID: AIQ_L10C5_FULL_0405
% Mức: VDC
Trình bày cách giải: Trước va chạm, hai xe chuyển động cùng chiều với động lượng $0,62$ và $0,11\ \mathrm{kg\,m/s}$. Tính tổng động lượng.
\loigiai{
Cùng chiều nên $p=p_1+p_2=0,62+0,11$. Suy ra kết quả $0,73\ \mathrm{kg\,m/s}$.
}
\end{bt}

% ===== Câu 82 =====
% ID: AIQ_L10C5_FULL_0406
% Mức: NB
\dangbt{Kiểm chứng định luật bảo toàn động lượng}
\begin{ex}
% ID: AIQ_L10C5_FULL_0406
% Mức: NB
Tổng động lượng trước va chạm là $0,5$ kg m/s, sau va chạm là $0,49$ kg m/s. Tính độ lệch tương đối so với trước.
\choice
{\True $2\ \mathrm{\%}$}
{$0,01\ \mathrm{\%}$}
{$4\ \mathrm{\%}$}
{$1\ \mathrm{\%}$}
\loigiai{
$\delta=|p_s-p_t|/p_t\cdot100\%=|0,49-0,5|/0,5\cdot100\%$. Suy ra kết quả $2\ \mathrm{\%}$.
}
\end{ex}

% ===== Câu 83 =====
% ID: AIQ_L10C5_FULL_0407
% Mức: NB
\begin{ex}
% ID: AIQ_L10C5_FULL_0407
% Mức: NB
Tổng động lượng trước va chạm là $0,52$ kg m/s, sau va chạm là $0,5$ kg m/s. Tính độ lệch tương đối so với trước.
\choice
{$0,02\ \mathrm{\%}$}
{\True $3,85\ \mathrm{\%}$}
{$7,69\ \mathrm{\%}$}
{$1,92\ \mathrm{\%}$}
\loigiai{
$\delta=|p_s-p_t|/p_t\cdot100\%=|0,5-0,52|/0,52\cdot100\%$. Suy ra kết quả $3,85\ \mathrm{\%}$.
}
\end{ex}

% ===== Câu 84 =====
% ID: AIQ_L10C5_FULL_0408
% Mức: TH
\begin{ex}
% ID: AIQ_L10C5_FULL_0408
% Mức: TH
Tổng động lượng trước va chạm là $0,54$ kg m/s, sau va chạm là $0,51$ kg m/s. Tính độ lệch tương đối so với trước.
\choice
{$0,03\ \mathrm{\%}$}
{$11,11\ \mathrm{\%}$}
{\True $5,56\ \mathrm{\%}$}
{$2,78\ \mathrm{\%}$}
\loigiai{
$\delta=|p_s-p_t|/p_t\cdot100\%=|0,51-0,54|/0,54\cdot100\%$. Suy ra kết quả $5,56\ \mathrm{\%}$.
}
\end{ex}

% ===== Câu 85 =====
% ID: AIQ_L10C5_FULL_0409
% Mức: TH
\begin{ex}
% ID: AIQ_L10C5_FULL_0409
% Mức: TH
Tổng động lượng trước va chạm là $0,56$ kg m/s, sau va chạm là $0,55$ kg m/s. Tính độ lệch tương đối so với trước.
\choice
{$0,01\ \mathrm{\%}$}
{$3,57\ \mathrm{\%}$}
{$0,89\ \mathrm{\%}$}
{\True $1,79\ \mathrm{\%}$}
\loigiai{
$\delta=|p_s-p_t|/p_t\cdot100\%=|0,55-0,56|/0,56\cdot100\%$. Suy ra kết quả $1,79\ \mathrm{\%}$.
}
\end{ex}

% ===== Câu 86 =====
% ID: AIQ_L10C5_FULL_0410
% Mức: VD
\begin{ex}
% ID: AIQ_L10C5_FULL_0410
% Mức: VD
Tổng động lượng trước va chạm là $0,58$ kg m/s, sau va chạm là $0,56$ kg m/s. Tính độ lệch tương đối so với trước.
\choice
{\True $3,45\ \mathrm{\%}$}
{$0,02\ \mathrm{\%}$}
{$6,9\ \mathrm{\%}$}
{$1,72\ \mathrm{\%}$}
\loigiai{
$\delta=|p_s-p_t|/p_t\cdot100\%=|0,56-0,58|/0,58\cdot100\%$. Suy ra kết quả $3,45\ \mathrm{\%}$.
}
\end{ex}

% ===== Câu 87 =====
% ID: AIQ_L10C5_FULL_0411
% Mức: VD
\begin{ex}
% ID: AIQ_L10C5_FULL_0411
% Mức: VD
Tổng động lượng trước va chạm là $0,6$ kg m/s, sau va chạm là $0,57$ kg m/s. Tính độ lệch tương đối so với trước.
\choice
{$0,03\ \mathrm{\%}$}
{\True $5\ \mathrm{\%}$}
{$10\ \mathrm{\%}$}
{$2,5\ \mathrm{\%}$}
\loigiai{
$\delta=|p_s-p_t|/p_t\cdot100\%=|0,57-0,6|/0,6\cdot100\%$. Suy ra kết quả $5\ \mathrm{\%}$.
}
\end{ex}

% ===== Câu 88 =====
% ID: AIQ_L10C5_FULL_0412
% Mức: VD
\begin{ex}
% ID: AIQ_L10C5_FULL_0412
% Mức: VD
Tổng động lượng trước va chạm là $0,62$ kg m/s, sau va chạm là $0,61$ kg m/s. Tính độ lệch tương đối so với trước.
\choice
{$0,01\ \mathrm{\%}$}
{$3,23\ \mathrm{\%}$}
{\True $1,61\ \mathrm{\%}$}
{$0,81\ \mathrm{\%}$}
\loigiai{
$\delta=|p_s-p_t|/p_t\cdot100\%=|0,61-0,62|/0,62\cdot100\%$. Suy ra kết quả $1,61\ \mathrm{\%}$.
}
\end{ex}

% ===== Câu 89 =====
% ID: AIQ_L10C5_FULL_0413
% Mức: VDC
\begin{ex}
% ID: AIQ_L10C5_FULL_0413
% Mức: VDC
Tổng động lượng trước va chạm là $0,64$ kg m/s, sau va chạm là $0,62$ kg m/s. Tính độ lệch tương đối so với trước.
\choice
{$0,02\ \mathrm{\%}$}
{$6,25\ \mathrm{\%}$}
{$1,56\ \mathrm{\%}$}
{\True $3,13\ \mathrm{\%}$}
\loigiai{
$\delta=|p_s-p_t|/p_t\cdot100\%=|0,62-0,64|/0,64\cdot100\%$. Suy ra kết quả $3,13\ \mathrm{\%}$.
}
\end{ex}

% ===== Câu 90 =====
% ID: AIQ_L10C5_FULL_0414
% Mức: VDC
\begin{ex}
% ID: AIQ_L10C5_FULL_0414
% Mức: VDC
Tổng động lượng trước va chạm là $0,66$ kg m/s, sau va chạm là $0,63$ kg m/s. Tính độ lệch tương đối so với trước.
\choice
{\True $4,55\ \mathrm{\%}$}
{$0,03\ \mathrm{\%}$}
{$9,09\ \mathrm{\%}$}
{$2,27\ \mathrm{\%}$}
\loigiai{
$\delta=|p_s-p_t|/p_t\cdot100\%=|0,63-0,66|/0,66\cdot100\%$. Suy ra kết quả $4,55\ \mathrm{\%}$.
}
\end{ex}

% ===== Câu 91 =====
% ID: AIQ_L10C5_FULL_0415
% Mức: VDC
\begin{ex}
% ID: AIQ_L10C5_FULL_0415
% Mức: VDC
Tổng động lượng trước va chạm là $0,68$ kg m/s, sau va chạm là $0,67$ kg m/s. Tính độ lệch tương đối so với trước.
\choice
{$0,01\ \mathrm{\%}$}
{\True $1,47\ \mathrm{\%}$}
{$2,94\ \mathrm{\%}$}
{$0,74\ \mathrm{\%}$}
\loigiai{
$\delta=|p_s-p_t|/p_t\cdot100\%=|0,67-0,68|/0,68\cdot100\%$. Suy ra kết quả $1,47\ \mathrm{\%}$.
}
\end{ex}

% ===== Câu 92 =====
% ID: AIQ_L10C5_FULL_0416
% Mức: NB
\begin{ex}
% ID: AIQ_L10C5_FULL_0416
% Mức: NB
Xét các phát biểu thuộc dạng “Kiểm chứng định luật bảo toàn động lượng” ở tình huống số 1.
\choiceTF
{\True So sánh tổng động lượng trước và sau va chạm.}
{\True Sai lệch nhỏ trong giới hạn đo phù hợp với định luật.}
{Thực nghiệm phải cho hai giá trị tuyệt đối bằng nhau.}
{\True Cần tính sai số tương đối.}
\loigiai{
\begin{itemchoice}
\item (Đúng) So sánh tổng động lượng trước và sau va chạm. (Vì): Phù hợp với định nghĩa hoặc hệ thức vật lí. b) Đúng: Phù hợp với định nghĩa hoặc hệ thức vật lí. c) Sai: Không phù hợp với định nghĩa hoặc hệ thức vật lí. d) Đúng: Phù hợp với định nghĩa hoặc hệ thức vật lí
\item (Đúng) Sai lệch nhỏ trong giới hạn đo phù hợp với định luật.
\item (Sai) Thực nghiệm phải cho hai giá trị tuyệt đối bằng nhau.
\item (Đúng) Cần tính sai số tương đối.
\end{itemchoice}
}
\end{ex}

% ===== Câu 93 =====
% ID: AIQ_L10C5_FULL_0417
% Mức: TH
\begin{ex}
% ID: AIQ_L10C5_FULL_0417
% Mức: TH
Xét các phát biểu thuộc dạng “Kiểm chứng định luật bảo toàn động lượng” ở tình huống số 2.
\choiceTF
{\True So sánh tổng động lượng trước và sau va chạm.}
{\True Sai lệch nhỏ trong giới hạn đo phù hợp với định luật.}
{Thực nghiệm phải cho hai giá trị tuyệt đối bằng nhau.}
{\True Cần tính sai số tương đối.}
\loigiai{
\begin{itemchoice}
\item (Đúng) So sánh tổng động lượng trước và sau va chạm. (Vì): Phù hợp với định nghĩa hoặc hệ thức vật lí. b) Đúng: Phù hợp với định nghĩa hoặc hệ thức vật lí. c) Sai: Không phù hợp với định nghĩa hoặc hệ thức vật lí. d) Đúng: Phù hợp với định nghĩa hoặc hệ thức vật lí
\item (Đúng) Sai lệch nhỏ trong giới hạn đo phù hợp với định luật.
\item (Sai) Thực nghiệm phải cho hai giá trị tuyệt đối bằng nhau.
\item (Đúng) Cần tính sai số tương đối.
\end{itemchoice}
}
\end{ex}

% ===== Câu 94 =====
% ID: AIQ_L10C5_FULL_0418
% Mức: TH
\begin{ex}
% ID: AIQ_L10C5_FULL_0418
% Mức: TH
Xét các phát biểu thuộc dạng “Kiểm chứng định luật bảo toàn động lượng” ở tình huống số 3.
\choiceTF
{\True So sánh tổng động lượng trước và sau va chạm.}
{\True Sai lệch nhỏ trong giới hạn đo phù hợp với định luật.}
{Thực nghiệm phải cho hai giá trị tuyệt đối bằng nhau.}
{\True Cần tính sai số tương đối.}
\loigiai{
\begin{itemchoice}
\item (Đúng) So sánh tổng động lượng trước và sau va chạm. (Vì): Phù hợp với định nghĩa hoặc hệ thức vật lí. b) Đúng: Phù hợp với định nghĩa hoặc hệ thức vật lí. c) Sai: Không phù hợp với định nghĩa hoặc hệ thức vật lí. d) Đúng: Phù hợp với định nghĩa hoặc hệ thức vật lí
\item (Đúng) Sai lệch nhỏ trong giới hạn đo phù hợp với định luật.
\item (Sai) Thực nghiệm phải cho hai giá trị tuyệt đối bằng nhau.
\item (Đúng) Cần tính sai số tương đối.
\end{itemchoice}
}
\end{ex}

% ===== Câu 95 =====
% ID: AIQ_L10C5_FULL_0419
% Mức: VD
\begin{ex}
% ID: AIQ_L10C5_FULL_0419
% Mức: VD
Xét các phát biểu thuộc dạng “Kiểm chứng định luật bảo toàn động lượng” ở tình huống số 4.
\choiceTF
{\True So sánh tổng động lượng trước và sau va chạm.}
{\True Sai lệch nhỏ trong giới hạn đo phù hợp với định luật.}
{Thực nghiệm phải cho hai giá trị tuyệt đối bằng nhau.}
{\True Cần tính sai số tương đối.}
\loigiai{
\begin{itemchoice}
\item (Đúng) So sánh tổng động lượng trước và sau va chạm. (Vì): Phù hợp với định nghĩa hoặc hệ thức vật lí. b) Đúng: Phù hợp với định nghĩa hoặc hệ thức vật lí. c) Sai: Không phù hợp với định nghĩa hoặc hệ thức vật lí. d) Đúng: Phù hợp với định nghĩa hoặc hệ thức vật lí
\item (Đúng) Sai lệch nhỏ trong giới hạn đo phù hợp với định luật.
\item (Sai) Thực nghiệm phải cho hai giá trị tuyệt đối bằng nhau.
\item (Đúng) Cần tính sai số tương đối.
\end{itemchoice}
}
\end{ex}

% ===== Câu 96 =====
% ID: AIQ_L10C5_FULL_0420
% Mức: VDC
\begin{ex}
% ID: AIQ_L10C5_FULL_0420
% Mức: VDC
Xét các phát biểu thuộc dạng “Kiểm chứng định luật bảo toàn động lượng” ở tình huống số 5.
\choiceTF
{\True So sánh tổng động lượng trước và sau va chạm.}
{\True Sai lệch nhỏ trong giới hạn đo phù hợp với định luật.}
{Thực nghiệm phải cho hai giá trị tuyệt đối bằng nhau.}
{\True Cần tính sai số tương đối.}
\loigiai{
\begin{itemchoice}
\item (Đúng) So sánh tổng động lượng trước và sau va chạm. (Vì): Phù hợp với định nghĩa hoặc hệ thức vật lí. b) Đúng: Phù hợp với định nghĩa hoặc hệ thức vật lí. c) Sai: Không phù hợp với định nghĩa hoặc hệ thức vật lí. d) Đúng: Phù hợp với định nghĩa hoặc hệ thức vật lí
\item (Đúng) Sai lệch nhỏ trong giới hạn đo phù hợp với định luật.
\item (Sai) Thực nghiệm phải cho hai giá trị tuyệt đối bằng nhau.
\item (Đúng) Cần tính sai số tương đối.
\end{itemchoice}
}
\end{ex}

% ===== Câu 97 =====
% ID: AIQ_L10C5_FULL_0421
% Mức: TH
\begin{ex}
% ID: AIQ_L10C5_FULL_0421
% Mức: TH
Tổng động lượng trước va chạm là $0,7$ kg m/s, sau va chạm là $0,68$ kg m/s. Tính độ lệch tương đối so với trước.
\shortans{2,86}
\loigiai{
$\delta=|p_s-p_t|/p_t\cdot100\%=|0,68-0,7|/0,7\cdot100\%$. Suy ra kết quả $2,86\ \mathrm{\%}$. Đáp án trả lời ngắn không ghi đơn vị.
}
\end{ex}

% ===== Câu 98 =====
% ID: AIQ_L10C5_FULL_0422
% Mức: TH
\begin{ex}
% ID: AIQ_L10C5_FULL_0422
% Mức: TH
Tổng động lượng trước va chạm là $0,72$ kg m/s, sau va chạm là $0,69$ kg m/s. Tính độ lệch tương đối so với trước.
\shortans{4,17}
\loigiai{
$\delta=|p_s-p_t|/p_t\cdot100\%=|0,69-0,72|/0,72\cdot100\%$. Suy ra kết quả $4,17\ \mathrm{\%}$. Đáp án trả lời ngắn không ghi đơn vị.
}
\end{ex}

% ===== Câu 99 =====
% ID: AIQ_L10C5_FULL_0423
% Mức: VD
\begin{ex}
% ID: AIQ_L10C5_FULL_0423
% Mức: VD
Tổng động lượng trước va chạm là $0,74$ kg m/s, sau va chạm là $0,73$ kg m/s. Tính độ lệch tương đối so với trước.
\shortans{1,35}
\loigiai{
$\delta=|p_s-p_t|/p_t\cdot100\%=|0,73-0,74|/0,74\cdot100\%$. Suy ra kết quả $1,35\ \mathrm{\%}$. Đáp án trả lời ngắn không ghi đơn vị.
}
\end{ex}

% ===== Câu 100 =====
% ID: AIQ_L10C5_FULL_0424
% Mức: VD
\begin{ex}
% ID: AIQ_L10C5_FULL_0424
% Mức: VD
Tổng động lượng trước va chạm là $0,76$ kg m/s, sau va chạm là $0,74$ kg m/s. Tính độ lệch tương đối so với trước.
\shortans{2,63}
\loigiai{
$\delta=|p_s-p_t|/p_t\cdot100\%=|0,74-0,76|/0,76\cdot100\%$. Suy ra kết quả $2,63\ \mathrm{\%}$. Đáp án trả lời ngắn không ghi đơn vị.
}
\end{ex}

% ===== Câu 101 =====
% ID: AIQ_L10C5_FULL_0425
% Mức: VDC
\begin{ex}
% ID: AIQ_L10C5_FULL_0425
% Mức: VDC
Tổng động lượng trước va chạm là $0,78$ kg m/s, sau va chạm là $0,75$ kg m/s. Tính độ lệch tương đối so với trước.
\shortans{3,85}
\loigiai{
$\delta=|p_s-p_t|/p_t\cdot100\%=|0,75-0,78|/0,78\cdot100\%$. Suy ra kết quả $3,85\ \mathrm{\%}$. Đáp án trả lời ngắn không ghi đơn vị.
}
\end{ex}

% ===== Câu 102 =====
% ID: AIQ_L10C5_FULL_0426
% Mức: VDC
\begin{ex}
% ID: AIQ_L10C5_FULL_0426
% Mức: VDC
Tổng động lượng trước va chạm là $0,8$ kg m/s, sau va chạm là $0,79$ kg m/s. Tính độ lệch tương đối so với trước.
\shortans{1,25}
\loigiai{
$\delta=|p_s-p_t|/p_t\cdot100\%=|0,79-0,8|/0,8\cdot100\%$. Suy ra kết quả $1,25\ \mathrm{\%}$. Đáp án trả lời ngắn không ghi đơn vị.
}
\end{ex}

% ===== Câu 103 =====
% ID: AIQ_L10C5_FULL_0427
% Mức: TH
\begin{bt}
% ID: AIQ_L10C5_FULL_0427
% Mức: TH
Trình bày cách giải: Tổng động lượng trước va chạm là $0,82$ kg m/s, sau va chạm là $0,8$ kg m/s. Tính độ lệch tương đối so với trước.
\loigiai{
$\delta=|p_s-p_t|/p_t\cdot100\%=|0,8-0,82|/0,82\cdot100\%$. Suy ra kết quả $2,44\ \mathrm{\%}$.
}
\end{bt}

% ===== Câu 104 =====
% ID: AIQ_L10C5_FULL_0428
% Mức: TH
\begin{bt}
% ID: AIQ_L10C5_FULL_0428
% Mức: TH
Trình bày cách giải: Tổng động lượng trước va chạm là $0,84$ kg m/s, sau va chạm là $0,81$ kg m/s. Tính độ lệch tương đối so với trước.
\loigiai{
$\delta=|p_s-p_t|/p_t\cdot100\%=|0,81-0,84|/0,84\cdot100\%$. Suy ra kết quả $3,57\ \mathrm{\%}$.
}
\end{bt}

% ===== Câu 105 =====
% ID: AIQ_L10C5_FULL_0429
% Mức: VD
\begin{bt}
% ID: AIQ_L10C5_FULL_0429
% Mức: VD
Trình bày cách giải: Tổng động lượng trước va chạm là $0,86$ kg m/s, sau va chạm là $0,85$ kg m/s. Tính độ lệch tương đối so với trước.
\loigiai{
$\delta=|p_s-p_t|/p_t\cdot100\%=|0,85-0,86|/0,86\cdot100\%$. Suy ra kết quả $1,16\ \mathrm{\%}$.
}
\end{bt}

% ===== Câu 106 =====
% ID: AIQ_L10C5_FULL_0430
% Mức: VD
\begin{bt}
% ID: AIQ_L10C5_FULL_0430
% Mức: VD
Trình bày cách giải: Tổng động lượng trước va chạm là $0,88$ kg m/s, sau va chạm là $0,86$ kg m/s. Tính độ lệch tương đối so với trước.
\loigiai{
$\delta=|p_s-p_t|/p_t\cdot100\%=|0,86-0,88|/0,88\cdot100\%$. Suy ra kết quả $2,27\ \mathrm{\%}$.
}
\end{bt}

% ===== Câu 107 =====
% ID: AIQ_L10C5_FULL_0431
% Mức: VDC
\begin{bt}
% ID: AIQ_L10C5_FULL_0431
% Mức: VDC
Trình bày cách giải: Tổng động lượng trước va chạm là $0,9$ kg m/s, sau va chạm là $0,87$ kg m/s. Tính độ lệch tương đối so với trước.
\loigiai{
$\delta=|p_s-p_t|/p_t\cdot100\%=|0,87-0,9|/0,9\cdot100\%$. Suy ra kết quả $3,33\ \mathrm{\%}$.
}
\end{bt}

% ===== Câu 108 =====
% ID: AIQ_L10C5_FULL_0432
% Mức: VDC
\begin{bt}
% ID: AIQ_L10C5_FULL_0432
% Mức: VDC
Trình bày cách giải: Tổng động lượng trước va chạm là $0,92$ kg m/s, sau va chạm là $0,91$ kg m/s. Tính độ lệch tương đối so với trước.
\loigiai{
$\delta=|p_s-p_t|/p_t\cdot100\%=|0,91-0,92|/0,92\cdot100\%$. Suy ra kết quả $1,09\ \mathrm{\%}$.
}
\end{bt}

% ===== Câu 109 =====
% ID: AIQ_L10C5_FULL_0433
% Mức: NB
\dangbt{Đánh giá sai số và nguyên nhân sai lệch}
\begin{ex}
% ID: AIQ_L10C5_FULL_0433
% Mức: NB
Kết quả thí nghiệm cho $p_{trước}=0,4$ và $p_{sau}=0,38\ \mathrm{kg\,m/s}$. Tính sai số tương đối.
\choice
{\True $5\ \mathrm{\%}$}
{$2\ \mathrm{\%}$}
{$10\ \mathrm{\%}$}
{$2,5\ \mathrm{\%}$}
\loigiai{
$\delta=|p_{sau}-p_{trước}|/p_{trước}\cdot100\%$. Suy ra kết quả $5\ \mathrm{\%}$.
}
\end{ex}

% ===== Câu 110 =====
% ID: AIQ_L10C5_FULL_0434
% Mức: NB
\begin{ex}
% ID: AIQ_L10C5_FULL_0434
% Mức: NB
Kết quả thí nghiệm cho $p_{trước}=0,43$ và $p_{sau}=0,41\ \mathrm{kg\,m/s}$. Tính sai số tương đối.
\choice
{$2\ \mathrm{\%}$}
{\True $4,65\ \mathrm{\%}$}
{$9,3\ \mathrm{\%}$}
{$2,33\ \mathrm{\%}$}
\loigiai{
$\delta=|p_{sau}-p_{trước}|/p_{trước}\cdot100\%$. Suy ra kết quả $4,65\ \mathrm{\%}$.
}
\end{ex}

% ===== Câu 111 =====
% ID: AIQ_L10C5_FULL_0435
% Mức: TH
\begin{ex}
% ID: AIQ_L10C5_FULL_0435
% Mức: TH
Kết quả thí nghiệm cho $p_{trước}=0,46$ và $p_{sau}=0,44\ \mathrm{kg\,m/s}$. Tính sai số tương đối.
\choice
{$2\ \mathrm{\%}$}
{$8,7\ \mathrm{\%}$}
{\True $4,35\ \mathrm{\%}$}
{$2,17\ \mathrm{\%}$}
\loigiai{
$\delta=|p_{sau}-p_{trước}|/p_{trước}\cdot100\%$. Suy ra kết quả $4,35\ \mathrm{\%}$.
}
\end{ex}

% ===== Câu 112 =====
% ID: AIQ_L10C5_FULL_0436
% Mức: TH
\begin{ex}
% ID: AIQ_L10C5_FULL_0436
% Mức: TH
Kết quả thí nghiệm cho $p_{trước}=0,49$ và $p_{sau}=0,47\ \mathrm{kg\,m/s}$. Tính sai số tương đối.
\choice
{$2\ \mathrm{\%}$}
{$8,16\ \mathrm{\%}$}
{$2,04\ \mathrm{\%}$}
{\True $4,08\ \mathrm{\%}$}
\loigiai{
$\delta=|p_{sau}-p_{trước}|/p_{trước}\cdot100\%$. Suy ra kết quả $4,08\ \mathrm{\%}$.
}
\end{ex}

% ===== Câu 113 =====
% ID: AIQ_L10C5_FULL_0437
% Mức: VD
\begin{ex}
% ID: AIQ_L10C5_FULL_0437
% Mức: VD
Kết quả thí nghiệm cho $p_{trước}=0,52$ và $p_{sau}=0,5\ \mathrm{kg\,m/s}$. Tính sai số tương đối.
\choice
{\True $3,85\ \mathrm{\%}$}
{$2\ \mathrm{\%}$}
{$7,69\ \mathrm{\%}$}
{$1,92\ \mathrm{\%}$}
\loigiai{
$\delta=|p_{sau}-p_{trước}|/p_{trước}\cdot100\%$. Suy ra kết quả $3,85\ \mathrm{\%}$.
}
\end{ex}

% ===== Câu 114 =====
% ID: AIQ_L10C5_FULL_0438
% Mức: VD
\begin{ex}
% ID: AIQ_L10C5_FULL_0438
% Mức: VD
Kết quả thí nghiệm cho $p_{trước}=0,55$ và $p_{sau}=0,53\ \mathrm{kg\,m/s}$. Tính sai số tương đối.
\choice
{$2\ \mathrm{\%}$}
{\True $3,64\ \mathrm{\%}$}
{$7,27\ \mathrm{\%}$}
{$1,82\ \mathrm{\%}$}
\loigiai{
$\delta=|p_{sau}-p_{trước}|/p_{trước}\cdot100\%$. Suy ra kết quả $3,64\ \mathrm{\%}$.
}
\end{ex}

% ===== Câu 115 =====
% ID: AIQ_L10C5_FULL_0439
% Mức: VD
\begin{ex}
% ID: AIQ_L10C5_FULL_0439
% Mức: VD
Kết quả thí nghiệm cho $p_{trước}=0,58$ và $p_{sau}=0,56\ \mathrm{kg\,m/s}$. Tính sai số tương đối.
\choice
{$2\ \mathrm{\%}$}
{$6,9\ \mathrm{\%}$}
{\True $3,45\ \mathrm{\%}$}
{$1,72\ \mathrm{\%}$}
\loigiai{
$\delta=|p_{sau}-p_{trước}|/p_{trước}\cdot100\%$. Suy ra kết quả $3,45\ \mathrm{\%}$.
}
\end{ex}

% ===== Câu 116 =====
% ID: AIQ_L10C5_FULL_0440
% Mức: VDC
\begin{ex}
% ID: AIQ_L10C5_FULL_0440
% Mức: VDC
Kết quả thí nghiệm cho $p_{trước}=0,61$ và $p_{sau}=0,59\ \mathrm{kg\,m/s}$. Tính sai số tương đối.
\choice
{$2\ \mathrm{\%}$}
{$6,56\ \mathrm{\%}$}
{$1,64\ \mathrm{\%}$}
{\True $3,28\ \mathrm{\%}$}
\loigiai{
$\delta=|p_{sau}-p_{trước}|/p_{trước}\cdot100\%$. Suy ra kết quả $3,28\ \mathrm{\%}$.
}
\end{ex}

% ===== Câu 117 =====
% ID: AIQ_L10C5_FULL_0441
% Mức: VDC
\begin{ex}
% ID: AIQ_L10C5_FULL_0441
% Mức: VDC
Kết quả thí nghiệm cho $p_{trước}=0,64$ và $p_{sau}=0,62\ \mathrm{kg\,m/s}$. Tính sai số tương đối.
\choice
{\True $3,13\ \mathrm{\%}$}
{$2\ \mathrm{\%}$}
{$6,25\ \mathrm{\%}$}
{$1,56\ \mathrm{\%}$}
\loigiai{
$\delta=|p_{sau}-p_{trước}|/p_{trước}\cdot100\%$. Suy ra kết quả $3,13\ \mathrm{\%}$.
}
\end{ex}

% ===== Câu 118 =====
% ID: AIQ_L10C5_FULL_0442
% Mức: VDC
\begin{ex}
% ID: AIQ_L10C5_FULL_0442
% Mức: VDC
Kết quả thí nghiệm cho $p_{trước}=0,67$ và $p_{sau}=0,65\ \mathrm{kg\,m/s}$. Tính sai số tương đối.
\choice
{$2\ \mathrm{\%}$}
{\True $2,99\ \mathrm{\%}$}
{$5,97\ \mathrm{\%}$}
{$1,49\ \mathrm{\%}$}
\loigiai{
$\delta=|p_{sau}-p_{trước}|/p_{trước}\cdot100\%$. Suy ra kết quả $2,99\ \mathrm{\%}$.
}
\end{ex}

% ===== Câu 119 =====
% ID: AIQ_L10C5_FULL_0443
% Mức: NB
\begin{ex}
% ID: AIQ_L10C5_FULL_0443
% Mức: NB
Xét các phát biểu thuộc dạng “Đánh giá sai số và nguyên nhân sai lệch” ở tình huống số 1.
\choiceTF
{\True Ma sát làm sai lệch kết quả.}
{\True Cổng quang đặt lệch có thể gây sai số.}
{Đo lặp loại bỏ hoàn toàn sai số hệ thống.}
{\True Cần kiểm tra số 0 của dụng cụ.}
\loigiai{
\begin{itemchoice}
\item (Đúng) Ma sát làm sai lệch kết quả. (Vì): Phù hợp với định nghĩa hoặc hệ thức vật lí. b) Đúng: Phù hợp với định nghĩa hoặc hệ thức vật lí. c) Sai: Không phù hợp với định nghĩa hoặc hệ thức vật lí. d) Đúng: Phù hợp với định nghĩa hoặc hệ thức vật lí
\item (Đúng) Cổng quang đặt lệch có thể gây sai số.
\item (Sai) Đo lặp loại bỏ hoàn toàn sai số hệ thống.
\item (Đúng) Cần kiểm tra số 0 của dụng cụ.
\end{itemchoice}
}
\end{ex}

% ===== Câu 120 =====
% ID: AIQ_L10C5_FULL_0444
% Mức: TH
\begin{ex}
% ID: AIQ_L10C5_FULL_0444
% Mức: TH
Xét các phát biểu thuộc dạng “Đánh giá sai số và nguyên nhân sai lệch” ở tình huống số 2.
\choiceTF
{\True Ma sát làm sai lệch kết quả.}
{\True Cổng quang đặt lệch có thể gây sai số.}
{Đo lặp loại bỏ hoàn toàn sai số hệ thống.}
{\True Cần kiểm tra số 0 của dụng cụ.}
\loigiai{
\begin{itemchoice}
\item (Đúng) Ma sát làm sai lệch kết quả. (Vì): Phù hợp với định nghĩa hoặc hệ thức vật lí. b) Đúng: Phù hợp với định nghĩa hoặc hệ thức vật lí. c) Sai: Không phù hợp với định nghĩa hoặc hệ thức vật lí. d) Đúng: Phù hợp với định nghĩa hoặc hệ thức vật lí
\item (Đúng) Cổng quang đặt lệch có thể gây sai số.
\item (Sai) Đo lặp loại bỏ hoàn toàn sai số hệ thống.
\item (Đúng) Cần kiểm tra số 0 của dụng cụ.
\end{itemchoice}
}
\end{ex}

% ===== Câu 121 =====
% ID: AIQ_L10C5_FULL_0445
% Mức: TH
\begin{ex}
% ID: AIQ_L10C5_FULL_0445
% Mức: TH
Xét các phát biểu thuộc dạng “Đánh giá sai số và nguyên nhân sai lệch” ở tình huống số 3.
\choiceTF
{\True Ma sát làm sai lệch kết quả.}
{\True Cổng quang đặt lệch có thể gây sai số.}
{Đo lặp loại bỏ hoàn toàn sai số hệ thống.}
{\True Cần kiểm tra số 0 của dụng cụ.}
\loigiai{
\begin{itemchoice}
\item (Đúng) Ma sát làm sai lệch kết quả. (Vì): Phù hợp với định nghĩa hoặc hệ thức vật lí. b) Đúng: Phù hợp với định nghĩa hoặc hệ thức vật lí. c) Sai: Không phù hợp với định nghĩa hoặc hệ thức vật lí. d) Đúng: Phù hợp với định nghĩa hoặc hệ thức vật lí
\item (Đúng) Cổng quang đặt lệch có thể gây sai số.
\item (Sai) Đo lặp loại bỏ hoàn toàn sai số hệ thống.
\item (Đúng) Cần kiểm tra số 0 của dụng cụ.
\end{itemchoice}
}
\end{ex}

% ===== Câu 122 =====
% ID: AIQ_L10C5_FULL_0446
% Mức: VD
\begin{ex}
% ID: AIQ_L10C5_FULL_0446
% Mức: VD
Xét các phát biểu thuộc dạng “Đánh giá sai số và nguyên nhân sai lệch” ở tình huống số 4.
\choiceTF
{\True Ma sát làm sai lệch kết quả.}
{\True Cổng quang đặt lệch có thể gây sai số.}
{Đo lặp loại bỏ hoàn toàn sai số hệ thống.}
{\True Cần kiểm tra số 0 của dụng cụ.}
\loigiai{
\begin{itemchoice}
\item (Đúng) Ma sát làm sai lệch kết quả. (Vì): Phù hợp với định nghĩa hoặc hệ thức vật lí. b) Đúng: Phù hợp với định nghĩa hoặc hệ thức vật lí. c) Sai: Không phù hợp với định nghĩa hoặc hệ thức vật lí. d) Đúng: Phù hợp với định nghĩa hoặc hệ thức vật lí
\item (Đúng) Cổng quang đặt lệch có thể gây sai số.
\item (Sai) Đo lặp loại bỏ hoàn toàn sai số hệ thống.
\item (Đúng) Cần kiểm tra số 0 của dụng cụ.
\end{itemchoice}
}
\end{ex}

% ===== Câu 123 =====
% ID: AIQ_L10C5_FULL_0447
% Mức: VDC
\begin{ex}
% ID: AIQ_L10C5_FULL_0447
% Mức: VDC
Xét các phát biểu thuộc dạng “Đánh giá sai số và nguyên nhân sai lệch” ở tình huống số 5.
\choiceTF
{\True Ma sát làm sai lệch kết quả.}
{\True Cổng quang đặt lệch có thể gây sai số.}
{Đo lặp loại bỏ hoàn toàn sai số hệ thống.}
{\True Cần kiểm tra số 0 của dụng cụ.}
\loigiai{
\begin{itemchoice}
\item (Đúng) Ma sát làm sai lệch kết quả. (Vì): Phù hợp với định nghĩa hoặc hệ thức vật lí. b) Đúng: Phù hợp với định nghĩa hoặc hệ thức vật lí. c) Sai: Không phù hợp với định nghĩa hoặc hệ thức vật lí. d) Đúng: Phù hợp với định nghĩa hoặc hệ thức vật lí
\item (Đúng) Cổng quang đặt lệch có thể gây sai số.
\item (Sai) Đo lặp loại bỏ hoàn toàn sai số hệ thống.
\item (Đúng) Cần kiểm tra số 0 của dụng cụ.
\end{itemchoice}
}
\end{ex}

% ===== Câu 124 =====
% ID: AIQ_L10C5_FULL_0448
% Mức: TH
\begin{ex}
% ID: AIQ_L10C5_FULL_0448
% Mức: TH
Kết quả thí nghiệm cho $p_{trước}=0,7$ và $p_{sau}=0,68\ \mathrm{kg\,m/s}$. Tính sai số tương đối.
\shortans{2,86}
\loigiai{
$\delta=|p_{sau}-p_{trước}|/p_{trước}\cdot100\%$. Suy ra kết quả $2,86\ \mathrm{\%}$. Đáp án trả lời ngắn không ghi đơn vị.
}
\end{ex}

% ===== Câu 125 =====
% ID: AIQ_L10C5_FULL_0449
% Mức: TH
\begin{ex}
% ID: AIQ_L10C5_FULL_0449
% Mức: TH
Kết quả thí nghiệm cho $p_{trước}=0,73$ và $p_{sau}=0,71\ \mathrm{kg\,m/s}$. Tính sai số tương đối.
\shortans{2,74}
\loigiai{
$\delta=|p_{sau}-p_{trước}|/p_{trước}\cdot100\%$. Suy ra kết quả $2,74\ \mathrm{\%}$. Đáp án trả lời ngắn không ghi đơn vị.
}
\end{ex}

% ===== Câu 126 =====
% ID: AIQ_L10C5_FULL_0450
% Mức: VD
\begin{ex}
% ID: AIQ_L10C5_FULL_0450
% Mức: VD
Kết quả thí nghiệm cho $p_{trước}=0,76$ và $p_{sau}=0,74\ \mathrm{kg\,m/s}$. Tính sai số tương đối.
\shortans{2,63}
\loigiai{
$\delta=|p_{sau}-p_{trước}|/p_{trước}\cdot100\%$. Suy ra kết quả $2,63\ \mathrm{\%}$. Đáp án trả lời ngắn không ghi đơn vị.
}
\end{ex}

% ===== Câu 127 =====
% ID: AIQ_L10C5_FULL_0451
% Mức: VD
\begin{ex}
% ID: AIQ_L10C5_FULL_0451
% Mức: VD
Kết quả thí nghiệm cho $p_{trước}=0,79$ và $p_{sau}=0,77\ \mathrm{kg\,m/s}$. Tính sai số tương đối.
\shortans{2,53}
\loigiai{
$\delta=|p_{sau}-p_{trước}|/p_{trước}\cdot100\%$. Suy ra kết quả $2,53\ \mathrm{\%}$. Đáp án trả lời ngắn không ghi đơn vị.
}
\end{ex}

% ===== Câu 128 =====
% ID: AIQ_L10C5_FULL_0452
% Mức: VDC
\begin{ex}
% ID: AIQ_L10C5_FULL_0452
% Mức: VDC
Kết quả thí nghiệm cho $p_{trước}=0,82$ và $p_{sau}=0,8\ \mathrm{kg\,m/s}$. Tính sai số tương đối.
\shortans{2,44}
\loigiai{
$\delta=|p_{sau}-p_{trước}|/p_{trước}\cdot100\%$. Suy ra kết quả $2,44\ \mathrm{\%}$. Đáp án trả lời ngắn không ghi đơn vị.
}
\end{ex}

% ===== Câu 129 =====
% ID: AIQ_L10C5_FULL_0453
% Mức: VDC
\begin{ex}
% ID: AIQ_L10C5_FULL_0453
% Mức: VDC
Kết quả thí nghiệm cho $p_{trước}=0,85$ và $p_{sau}=0,83\ \mathrm{kg\,m/s}$. Tính sai số tương đối.
\shortans{2,35}
\loigiai{
$\delta=|p_{sau}-p_{trước}|/p_{trước}\cdot100\%$. Suy ra kết quả $2,35\ \mathrm{\%}$. Đáp án trả lời ngắn không ghi đơn vị.
}
\end{ex}

% ===== Câu 130 =====
% ID: AIQ_L10C5_FULL_0454
% Mức: TH
\begin{bt}
% ID: AIQ_L10C5_FULL_0454
% Mức: TH
Trình bày cách giải: Kết quả thí nghiệm cho $p_{trước}=0,88$ và $p_{sau}=0,86\ \mathrm{kg\,m/s}$. Tính sai số tương đối.
\loigiai{
$\delta=|p_{sau}-p_{trước}|/p_{trước}\cdot100\%$. Suy ra kết quả $2,27\ \mathrm{\%}$.
}
\end{bt}

% ===== Câu 131 =====
% ID: AIQ_L10C5_FULL_0455
% Mức: TH
\begin{bt}
% ID: AIQ_L10C5_FULL_0455
% Mức: TH
Trình bày cách giải: Kết quả thí nghiệm cho $p_{trước}=0,91$ và $p_{sau}=0,89\ \mathrm{kg\,m/s}$. Tính sai số tương đối.
\loigiai{
$\delta=|p_{sau}-p_{trước}|/p_{trước}\cdot100\%$. Suy ra kết quả $2,2\ \mathrm{\%}$.
}
\end{bt}

% ===== Câu 132 =====
% ID: AIQ_L10C5_FULL_0456
% Mức: VD
\begin{bt}
% ID: AIQ_L10C5_FULL_0456
% Mức: VD
Trình bày cách giải: Kết quả thí nghiệm cho $p_{trước}=0,94$ và $p_{sau}=0,92\ \mathrm{kg\,m/s}$. Tính sai số tương đối.
\loigiai{
$\delta=|p_{sau}-p_{trước}|/p_{trước}\cdot100\%$. Suy ra kết quả $2,13\ \mathrm{\%}$.
}
\end{bt}

% ===== Câu 133 =====
% ID: AIQ_L10C5_FULL_0457
% Mức: VD
\begin{bt}
% ID: AIQ_L10C5_FULL_0457
% Mức: VD
Trình bày cách giải: Kết quả thí nghiệm cho $p_{trước}=0,97$ và $p_{sau}=0,95\ \mathrm{kg\,m/s}$. Tính sai số tương đối.
\loigiai{
$\delta=|p_{sau}-p_{trước}|/p_{trước}\cdot100\%$. Suy ra kết quả $2,06\ \mathrm{\%}$.
}
\end{bt}

% ===== Câu 134 =====
% ID: AIQ_L10C5_FULL_0458
% Mức: VDC
\begin{bt}
% ID: AIQ_L10C5_FULL_0458
% Mức: VDC
Trình bày cách giải: Kết quả thí nghiệm cho $p_{trước}=1$ và $p_{sau}=0,98\ \mathrm{kg\,m/s}$. Tính sai số tương đối.
\loigiai{
$\delta=|p_{sau}-p_{trước}|/p_{trước}\cdot100\%$. Suy ra kết quả $2\ \mathrm{\%}$.
}
\end{bt}

% ===== Câu 135 =====
% ID: AIQ_L10C5_FULL_0459
% Mức: VDC
\begin{bt}
% ID: AIQ_L10C5_FULL_0459
% Mức: VDC
Trình bày cách giải: Kết quả thí nghiệm cho $p_{trước}=1,03$ và $p_{sau}=1,01\ \mathrm{kg\,m/s}$. Tính sai số tương đối.
\loigiai{
$\delta=|p_{sau}-p_{trước}|/p_{trước}\cdot100\%$. Suy ra kết quả $1,94\ \mathrm{\%}$.
}
\end{bt}

% ===== Câu 136 =====
% ID: AIQ_L10C5_FULL_0460
% Mức: NB
\dangbt{Thiết kế, cải tiến và báo cáo thí nghiệm va chạm}
\begin{ex}
% ID: AIQ_L10C5_FULL_0460
% Mức: NB
Một phương án khảo sát 3 cấu hình va chạm, mỗi cấu hình đo lặp 4 lần. Tổng số lượt đo là bao nhiêu?
\choice
{\True $12\ \mathrm{lượt}$}
{$7\ \mathrm{lượt}$}
{$6\ \mathrm{lượt}$}
{$16\ \mathrm{lượt}$}
\loigiai{
Tổng số lượt đo bằng $3\cdot4=12$. Suy ra kết quả $12\ \mathrm{lượt}$.
}
\end{ex}

% ===== Câu 137 =====
% ID: AIQ_L10C5_FULL_0461
% Mức: NB
\begin{ex}
% ID: AIQ_L10C5_FULL_0461
% Mức: NB
Một phương án khảo sát 4 cấu hình va chạm, mỗi cấu hình đo lặp 4 lần. Tổng số lượt đo là bao nhiêu?
\choice
{$8\ \mathrm{lượt}$}
{\True $16\ \mathrm{lượt}$}
{$8\ \mathrm{lượt}$}
{$20\ \mathrm{lượt}$}
\loigiai{
Tổng số lượt đo bằng $4\cdot4=16$. Suy ra kết quả $16\ \mathrm{lượt}$.
}
\end{ex}

% ===== Câu 138 =====
% ID: AIQ_L10C5_FULL_0462
% Mức: TH
\begin{ex}
% ID: AIQ_L10C5_FULL_0462
% Mức: TH
Một phương án khảo sát 5 cấu hình va chạm, mỗi cấu hình đo lặp 4 lần. Tổng số lượt đo là bao nhiêu?
\choice
{$9\ \mathrm{lượt}$}
{$10\ \mathrm{lượt}$}
{\True $20\ \mathrm{lượt}$}
{$24\ \mathrm{lượt}$}
\loigiai{
Tổng số lượt đo bằng $5\cdot4=20$. Suy ra kết quả $20\ \mathrm{lượt}$.
}
\end{ex}

% ===== Câu 139 =====
% ID: AIQ_L10C5_FULL_0463
% Mức: TH
\begin{ex}
% ID: AIQ_L10C5_FULL_0463
% Mức: TH
Một phương án khảo sát 6 cấu hình va chạm, mỗi cấu hình đo lặp 4 lần. Tổng số lượt đo là bao nhiêu?
\choice
{$10\ \mathrm{lượt}$}
{$12\ \mathrm{lượt}$}
{$28\ \mathrm{lượt}$}
{\True $24\ \mathrm{lượt}$}
\loigiai{
Tổng số lượt đo bằng $6\cdot4=24$. Suy ra kết quả $24\ \mathrm{lượt}$.
}
\end{ex}

% ===== Câu 140 =====
% ID: AIQ_L10C5_FULL_0464
% Mức: VD
\begin{ex}
% ID: AIQ_L10C5_FULL_0464
% Mức: VD
Một phương án khảo sát 7 cấu hình va chạm, mỗi cấu hình đo lặp 4 lần. Tổng số lượt đo là bao nhiêu?
\choice
{\True $28\ \mathrm{lượt}$}
{$11\ \mathrm{lượt}$}
{$14\ \mathrm{lượt}$}
{$32\ \mathrm{lượt}$}
\loigiai{
Tổng số lượt đo bằng $7\cdot4=28$. Suy ra kết quả $28\ \mathrm{lượt}$.
}
\end{ex}

% ===== Câu 141 =====
% ID: AIQ_L10C5_FULL_0465
% Mức: VD
\begin{ex}
% ID: AIQ_L10C5_FULL_0465
% Mức: VD
Một phương án khảo sát 3 cấu hình va chạm, mỗi cấu hình đo lặp 4 lần. Tổng số lượt đo là bao nhiêu?
\choice
{$7\ \mathrm{lượt}$}
{\True $12\ \mathrm{lượt}$}
{$6\ \mathrm{lượt}$}
{$16\ \mathrm{lượt}$}
\loigiai{
Tổng số lượt đo bằng $3\cdot4=12$. Suy ra kết quả $12\ \mathrm{lượt}$.
}
\end{ex}

% ===== Câu 142 =====
% ID: AIQ_L10C5_FULL_0466
% Mức: VD
\begin{ex}
% ID: AIQ_L10C5_FULL_0466
% Mức: VD
Một phương án khảo sát 4 cấu hình va chạm, mỗi cấu hình đo lặp 4 lần. Tổng số lượt đo là bao nhiêu?
\choice
{$8\ \mathrm{lượt}$}
{$8\ \mathrm{lượt}$}
{\True $16\ \mathrm{lượt}$}
{$20\ \mathrm{lượt}$}
\loigiai{
Tổng số lượt đo bằng $4\cdot4=16$. Suy ra kết quả $16\ \mathrm{lượt}$.
}
\end{ex}

% ===== Câu 143 =====
% ID: AIQ_L10C5_FULL_0467
% Mức: VDC
\begin{ex}
% ID: AIQ_L10C5_FULL_0467
% Mức: VDC
Một phương án khảo sát 5 cấu hình va chạm, mỗi cấu hình đo lặp 4 lần. Tổng số lượt đo là bao nhiêu?
\choice
{$9\ \mathrm{lượt}$}
{$10\ \mathrm{lượt}$}
{$24\ \mathrm{lượt}$}
{\True $20\ \mathrm{lượt}$}
\loigiai{
Tổng số lượt đo bằng $5\cdot4=20$. Suy ra kết quả $20\ \mathrm{lượt}$.
}
\end{ex}

% ===== Câu 144 =====
% ID: AIQ_L10C5_FULL_0468
% Mức: VDC
\begin{ex}
% ID: AIQ_L10C5_FULL_0468
% Mức: VDC
Một phương án khảo sát 6 cấu hình va chạm, mỗi cấu hình đo lặp 4 lần. Tổng số lượt đo là bao nhiêu?
\choice
{\True $24\ \mathrm{lượt}$}
{$10\ \mathrm{lượt}$}
{$12\ \mathrm{lượt}$}
{$28\ \mathrm{lượt}$}
\loigiai{
Tổng số lượt đo bằng $6\cdot4=24$. Suy ra kết quả $24\ \mathrm{lượt}$.
}
\end{ex}

% ===== Câu 145 =====
% ID: AIQ_L10C5_FULL_0469
% Mức: VDC
\begin{ex}
% ID: AIQ_L10C5_FULL_0469
% Mức: VDC
Một phương án khảo sát 7 cấu hình va chạm, mỗi cấu hình đo lặp 4 lần. Tổng số lượt đo là bao nhiêu?
\choice
{$11\ \mathrm{lượt}$}
{\True $28\ \mathrm{lượt}$}
{$14\ \mathrm{lượt}$}
{$32\ \mathrm{lượt}$}
\loigiai{
Tổng số lượt đo bằng $7\cdot4=28$. Suy ra kết quả $28\ \mathrm{lượt}$.
}
\end{ex}

% ===== Câu 146 =====
% ID: AIQ_L10C5_FULL_0470
% Mức: NB
\begin{ex}
% ID: AIQ_L10C5_FULL_0470
% Mức: NB
Xét các phát biểu thuộc dạng “Thiết kế, cải tiến và báo cáo thí nghiệm va chạm” ở tình huống số 1.
\choiceTF
{\True Báo cáo cần nêu mục tiêu và dụng cụ.}
{\True Phải ghi số liệu trung thực.}
{Có thể bỏ qua đơn vị trong bảng.}
{\True Nên đề xuất cải tiến giảm ma sát.}
\loigiai{
\begin{itemchoice}
\item (Đúng) Báo cáo cần nêu mục tiêu và dụng cụ. (Vì): Phù hợp với định nghĩa hoặc hệ thức vật lí. b) Đúng: Phù hợp với định nghĩa hoặc hệ thức vật lí. c) Sai: Không phù hợp với định nghĩa hoặc hệ thức vật lí. d) Đúng: Phù hợp với định nghĩa hoặc hệ thức vật lí
\item (Đúng) Phải ghi số liệu trung thực.
\item (Sai) Có thể bỏ qua đơn vị trong bảng.
\item (Đúng) Nên đề xuất cải tiến giảm ma sát.
\end{itemchoice}
}
\end{ex}

% ===== Câu 147 =====
% ID: AIQ_L10C5_FULL_0471
% Mức: TH
\begin{ex}
% ID: AIQ_L10C5_FULL_0471
% Mức: TH
Xét các phát biểu thuộc dạng “Thiết kế, cải tiến và báo cáo thí nghiệm va chạm” ở tình huống số 2.
\choiceTF
{\True Báo cáo cần nêu mục tiêu và dụng cụ.}
{\True Phải ghi số liệu trung thực.}
{Có thể bỏ qua đơn vị trong bảng.}
{\True Nên đề xuất cải tiến giảm ma sát.}
\loigiai{
\begin{itemchoice}
\item (Đúng) Báo cáo cần nêu mục tiêu và dụng cụ. (Vì): Phù hợp với định nghĩa hoặc hệ thức vật lí. b) Đúng: Phù hợp với định nghĩa hoặc hệ thức vật lí. c) Sai: Không phù hợp với định nghĩa hoặc hệ thức vật lí. d) Đúng: Phù hợp với định nghĩa hoặc hệ thức vật lí
\item (Đúng) Phải ghi số liệu trung thực.
\item (Sai) Có thể bỏ qua đơn vị trong bảng.
\item (Đúng) Nên đề xuất cải tiến giảm ma sát.
\end{itemchoice}
}
\end{ex}

% ===== Câu 148 =====
% ID: AIQ_L10C5_FULL_0472
% Mức: TH
\begin{ex}
% ID: AIQ_L10C5_FULL_0472
% Mức: TH
Xét các phát biểu thuộc dạng “Thiết kế, cải tiến và báo cáo thí nghiệm va chạm” ở tình huống số 3.
\choiceTF
{\True Báo cáo cần nêu mục tiêu và dụng cụ.}
{\True Phải ghi số liệu trung thực.}
{Có thể bỏ qua đơn vị trong bảng.}
{\True Nên đề xuất cải tiến giảm ma sát.}
\loigiai{
\begin{itemchoice}
\item (Đúng) Báo cáo cần nêu mục tiêu và dụng cụ. (Vì): Phù hợp với định nghĩa hoặc hệ thức vật lí. b) Đúng: Phù hợp với định nghĩa hoặc hệ thức vật lí. c) Sai: Không phù hợp với định nghĩa hoặc hệ thức vật lí. d) Đúng: Phù hợp với định nghĩa hoặc hệ thức vật lí
\item (Đúng) Phải ghi số liệu trung thực.
\item (Sai) Có thể bỏ qua đơn vị trong bảng.
\item (Đúng) Nên đề xuất cải tiến giảm ma sát.
\end{itemchoice}
}
\end{ex}

% ===== Câu 149 =====
% ID: AIQ_L10C5_FULL_0473
% Mức: VD
\begin{ex}
% ID: AIQ_L10C5_FULL_0473
% Mức: VD
Xét các phát biểu thuộc dạng “Thiết kế, cải tiến và báo cáo thí nghiệm va chạm” ở tình huống số 4.
\choiceTF
{\True Báo cáo cần nêu mục tiêu và dụng cụ.}
{\True Phải ghi số liệu trung thực.}
{Có thể bỏ qua đơn vị trong bảng.}
{\True Nên đề xuất cải tiến giảm ma sát.}
\loigiai{
\begin{itemchoice}
\item (Đúng) Báo cáo cần nêu mục tiêu và dụng cụ. (Vì): Phù hợp với định nghĩa hoặc hệ thức vật lí. b) Đúng: Phù hợp với định nghĩa hoặc hệ thức vật lí. c) Sai: Không phù hợp với định nghĩa hoặc hệ thức vật lí. d) Đúng: Phù hợp với định nghĩa hoặc hệ thức vật lí
\item (Đúng) Phải ghi số liệu trung thực.
\item (Sai) Có thể bỏ qua đơn vị trong bảng.
\item (Đúng) Nên đề xuất cải tiến giảm ma sát.
\end{itemchoice}
}
\end{ex}

% ===== Câu 150 =====
% ID: AIQ_L10C5_FULL_0474
% Mức: VDC
\begin{ex}
% ID: AIQ_L10C5_FULL_0474
% Mức: VDC
Xét các phát biểu thuộc dạng “Thiết kế, cải tiến và báo cáo thí nghiệm va chạm” ở tình huống số 5.
\choiceTF
{\True Báo cáo cần nêu mục tiêu và dụng cụ.}
{\True Phải ghi số liệu trung thực.}
{Có thể bỏ qua đơn vị trong bảng.}
{\True Nên đề xuất cải tiến giảm ma sát.}
\loigiai{
\begin{itemchoice}
\item (Đúng) Báo cáo cần nêu mục tiêu và dụng cụ. (Vì): Phù hợp với định nghĩa hoặc hệ thức vật lí. b) Đúng: Phù hợp với định nghĩa hoặc hệ thức vật lí. c) Sai: Không phù hợp với định nghĩa hoặc hệ thức vật lí. d) Đúng: Phù hợp với định nghĩa hoặc hệ thức vật lí
\item (Đúng) Phải ghi số liệu trung thực.
\item (Sai) Có thể bỏ qua đơn vị trong bảng.
\item (Đúng) Nên đề xuất cải tiến giảm ma sát.
\end{itemchoice}
}
\end{ex}

% ===== Câu 151 =====
% ID: AIQ_L10C5_FULL_0475
% Mức: TH
\begin{ex}
% ID: AIQ_L10C5_FULL_0475
% Mức: TH
Một phương án khảo sát 3 cấu hình va chạm, mỗi cấu hình đo lặp 4 lần. Tổng số lượt đo là bao nhiêu?
\shortans{12}
\loigiai{
Tổng số lượt đo bằng $3\cdot4=12$. Suy ra kết quả $12\ \mathrm{lượt}$. Đáp án trả lời ngắn không ghi đơn vị.
}
\end{ex}

% ===== Câu 152 =====
% ID: AIQ_L10C5_FULL_0476
% Mức: TH
\begin{ex}
% ID: AIQ_L10C5_FULL_0476
% Mức: TH
Một phương án khảo sát 4 cấu hình va chạm, mỗi cấu hình đo lặp 4 lần. Tổng số lượt đo là bao nhiêu?
\shortans{16}
\loigiai{
Tổng số lượt đo bằng $4\cdot4=16$. Suy ra kết quả $16\ \mathrm{lượt}$. Đáp án trả lời ngắn không ghi đơn vị.
}
\end{ex}

% ===== Câu 153 =====
% ID: AIQ_L10C5_FULL_0477
% Mức: VD
\begin{ex}
% ID: AIQ_L10C5_FULL_0477
% Mức: VD
Một phương án khảo sát 5 cấu hình va chạm, mỗi cấu hình đo lặp 4 lần. Tổng số lượt đo là bao nhiêu?
\shortans{20}
\loigiai{
Tổng số lượt đo bằng $5\cdot4=20$. Suy ra kết quả $20\ \mathrm{lượt}$. Đáp án trả lời ngắn không ghi đơn vị.
}
\end{ex}

% ===== Câu 154 =====
% ID: AIQ_L10C5_FULL_0478
% Mức: VD
\begin{ex}
% ID: AIQ_L10C5_FULL_0478
% Mức: VD
Một phương án khảo sát 6 cấu hình va chạm, mỗi cấu hình đo lặp 4 lần. Tổng số lượt đo là bao nhiêu?
\shortans{24}
\loigiai{
Tổng số lượt đo bằng $6\cdot4=24$. Suy ra kết quả $24\ \mathrm{lượt}$. Đáp án trả lời ngắn không ghi đơn vị.
}
\end{ex}

% ===== Câu 155 =====
% ID: AIQ_L10C5_FULL_0479
% Mức: VDC
\begin{ex}
% ID: AIQ_L10C5_FULL_0479
% Mức: VDC
Một phương án khảo sát 7 cấu hình va chạm, mỗi cấu hình đo lặp 4 lần. Tổng số lượt đo là bao nhiêu?
\shortans{28}
\loigiai{
Tổng số lượt đo bằng $7\cdot4=28$. Suy ra kết quả $28\ \mathrm{lượt}$. Đáp án trả lời ngắn không ghi đơn vị.
}
\end{ex}

% ===== Câu 156 =====
% ID: AIQ_L10C5_FULL_0480
% Mức: VDC
\begin{ex}
% ID: AIQ_L10C5_FULL_0480
% Mức: VDC
Một phương án khảo sát 3 cấu hình va chạm, mỗi cấu hình đo lặp 4 lần. Tổng số lượt đo là bao nhiêu?
\shortans{12}
\loigiai{
Tổng số lượt đo bằng $3\cdot4=12$. Suy ra kết quả $12\ \mathrm{lượt}$. Đáp án trả lời ngắn không ghi đơn vị.
}
\end{ex}

% ===== Câu 157 =====
% ID: AIQ_L10C5_FULL_0481
% Mức: TH
\begin{bt}
% ID: AIQ_L10C5_FULL_0481
% Mức: TH
Trình bày cách giải: Một phương án khảo sát 4 cấu hình va chạm, mỗi cấu hình đo lặp 4 lần. Tổng số lượt đo là bao nhiêu?
\loigiai{
Tổng số lượt đo bằng $4\cdot4=16$. Suy ra kết quả $16\ \mathrm{lượt}$.
}
\end{bt}

% ===== Câu 158 =====
% ID: AIQ_L10C5_FULL_0482
% Mức: TH
\begin{bt}
% ID: AIQ_L10C5_FULL_0482
% Mức: TH
Trình bày cách giải: Một phương án khảo sát 5 cấu hình va chạm, mỗi cấu hình đo lặp 4 lần. Tổng số lượt đo là bao nhiêu?
\loigiai{
Tổng số lượt đo bằng $5\cdot4=20$. Suy ra kết quả $20\ \mathrm{lượt}$.
}
\end{bt}

% ===== Câu 159 =====
% ID: AIQ_L10C5_FULL_0483
% Mức: VD
\begin{bt}
% ID: AIQ_L10C5_FULL_0483
% Mức: VD
Trình bày cách giải: Một phương án khảo sát 6 cấu hình va chạm, mỗi cấu hình đo lặp 4 lần. Tổng số lượt đo là bao nhiêu?
\loigiai{
Tổng số lượt đo bằng $6\cdot4=24$. Suy ra kết quả $24\ \mathrm{lượt}$.
}
\end{bt}

% ===== Câu 160 =====
% ID: AIQ_L10C5_FULL_0484
% Mức: VD
\begin{bt}
% ID: AIQ_L10C5_FULL_0484
% Mức: VD
Trình bày cách giải: Một phương án khảo sát 7 cấu hình va chạm, mỗi cấu hình đo lặp 4 lần. Tổng số lượt đo là bao nhiêu?
\loigiai{
Tổng số lượt đo bằng $7\cdot4=28$. Suy ra kết quả $28\ \mathrm{lượt}$.
}
\end{bt}

% ===== Câu 161 =====
% ID: AIQ_L10C5_FULL_0485
% Mức: VDC
\begin{bt}
% ID: AIQ_L10C5_FULL_0485
% Mức: VDC
Trình bày cách giải: Một phương án khảo sát 3 cấu hình va chạm, mỗi cấu hình đo lặp 4 lần. Tổng số lượt đo là bao nhiêu?
\loigiai{
Tổng số lượt đo bằng $3\cdot4=12$. Suy ra kết quả $12\ \mathrm{lượt}$.
}
\end{bt}

% ===== Câu 162 =====
% ID: AIQ_L10C5_FULL_0486
% Mức: VDC
\begin{bt}
% ID: AIQ_L10C5_FULL_0486
% Mức: VDC
Trình bày cách giải: Một phương án khảo sát 4 cấu hình va chạm, mỗi cấu hình đo lặp 4 lần. Tổng số lượt đo là bao nhiêu?
\loigiai{
Tổng số lượt đo bằng $4\cdot4=16$. Suy ra kết quả $16\ \mathrm{lượt}$.
}
\end{bt}
